% Appendix A (continued), Parts C and D: manifolds, germs and cones; groups, symmetry and conservation.
% Source: site/sections/13-appendix.html (#part-c, #part-d), converted in full; entry numbers C.1--C.22 and D.1--D.34
% are the site's. Claims about numbered results follow theory/propositions.json (public_statement).
% No \section here: this file continues Appendix A (app:background), whose Parts A and B are in A1-background.tex.
% "Used in" lines use the macro of A1-background.tex; this copy (identical) lets the file compile on its own.
\providecommand{\fbAusedin}[1]{\par\smallskip\noindent{\small\color{inkgray}{\sffamily\bfseries\color{tide}Used in.}\enspace #1\par}\medskip}

% ======================================================================================================================
\subsection*{Part C. Manifolds, germs and cones}
\phantomsection\label{app:bg-c}\addcontentsline{toc}{subsection}{Part C. Manifolds, germs and cones}

Parts A and B used only the arrows of a method and its wiring. Part C turns to its geometry. A PEFT method in this paper
is a smooth map $\rho:(Q,q_0)\to(\Theta,\theta_0)$, and its expressivity is read off the set $\Im M=\rho(Q)$ near
$\theta_0$. This part collects the geometry behind that reading, from manifolds, differentials and rank for the map to
germs, cones, singular points and semialgebraic dimension for the image, which is often not a manifold.

% ----------------------------------------------------------------------------------------------------------------------
\subsubsection*{Manifolds and smooth maps}

\begin{bgentry}{C.1}{smooth manifolds and smooth maps}
A \textbf{smooth manifold} of dimension $k$ is a space covered by \textbf{charts}, bijections from open pieces of it
onto open subsets of $\R^k$, whose changes of chart are smooth. Every manifold in this paper is an open subset of a
vector space or an embedded submanifold of one (\bgref{C.4}), or is diffeomorphic to one, as quotients (\bgref{D.22})
and the Grassmannian (\bgref{C.15}) are. A \textbf{diffeomorphism} is a smooth bijection with smooth inverse, and
$\mathrm{Diff}(Q)$ is the group of diffeomorphisms of $Q$ under composition. A \textbf{local diffeomorphism} maps an
open neighbourhood of each point diffeomorphically onto an open set; by the inverse function theorem, $\rho$ does this
near $q$ whenever $d\rho_q$ (\bgref{C.2}) is invertible.
\end{bgentry}

LoRA's \citep{hu2021lora} parameter space $\R^{m\times r}\times\R^{r\times n}$ is a vector space, and HRA's
\citep{yuan2024bridging} $\{(u_1,\dots,u_r):u_i\ne0\}$ is open in $(\R^n)^r$. The rotation group $\SO(n)$ is a compact
manifold of dimension $n(n-1)/2$, and the Cayley map (\bgref{D.7}) is a diffeomorphism from $\so(n)$ onto an open dense
subset of it. A method is a pointed smooth map, and its gauge group
$\{\varphi\in\mathrm{Diff}(Q):\rho\circ\varphi=\rho\}$ (\bgref{D.21}) contains LoRA's $(B,A)\mapsto(Bg^{-1},gA)$ for
every $g\in\GL_r$.

\fbAusedin{\thmref{Definition}{I.2}, \thmref{Definition}{I.5}, \thmref{Theorem}{II.7}, \thmref{Proposition}{IV.3}.
  \textsf{Reading:} \citet[Ch.~1 and~2]{xlee2013introduction}.}

\begin{bgentry}{C.2}{tangent spaces, differentials and covectors}
For a manifold $Q\subseteq\R^N$, the \textbf{tangent space} $T_qQ$ is the set of velocities $c'(0)$ of smooth curves $c$
in $Q$ with $c(0)=q$. It is a linear space of dimension $\dim Q$, all of $\R^N$ when $Q$ is open. The
\textbf{differential} of a smooth $\rho$ at $q$ is the linear map $d\rho_q:T_qQ\to T_{\rho(q)}\Theta$,
$c'(0)\mapsto(\rho\circ c)'(0)$. In coordinates it is the Jacobian, and the chain rule reads $d(\rho\circ
h)_q=d\rho_{h(q)}\circ dh_q$. The \textbf{cotangent space} $T^*_\theta\Theta$ is the dual of $T_\theta\Theta$. The
differential $dL_\theta$ of a loss is a \textbf{covector}, and it becomes a gradient vector $\nabla L(\theta)$ only once
an inner product is chosen.
\end{bgentry}

For LoRA, $d\rho_{(B,A)}(\dot B,\dot A)=s(\dot BA+B\dot A)$. For a rotation orbit, $T_I\SO(n)=\so(n)$ and
$T_{W_0R}\big(W_0\cdot\SO(n)\big)=W_0R\,\so(n)$. \thmref{Observation}{I.6} is the chain rule at $q_0$.
\thmref{Theorem}{III.5} drives the parameters by any cotangent signal $\gamma(t)\in T^*_{\rho(q)}\Theta$, pulled back to
$Q$ as $\gamma\circ d\rho_q$, which is $D\rho\T\gamma$ in coordinates.

\fbAusedin{the notation (Table~\ref{tab:notation}), \thmref{Definition}{I.5}, \thmref{Observation}{I.6},
  \thmref{Theorem}{II.2}, \thmref{Theorem}{II.7}, \thmref{Theorem}{III.5}, \thmref{Definition}{IV.1}. \textsf{Reading:}
  \citet[Ch.~3 and~11]{xlee2013introduction}.}

\begin{bgentry}{C.3}{rank, expressive dimension and first-order image}
The \textbf{rank} of $\rho$ at $q$ is $\rk d\rho_q$. It is lower semicontinuous, since a nonzero $k\times k$ minor of
the Jacobian stays nonzero nearby, so $\{q:\rk d\rho_q\ge k\}$ is open. The \textbf{generic rank} is $\max_q\rk
d\rho_q$. If $\rho$ is real-analytic and $Q$ is connected, the generic rank is attained on an open dense set whose
complement has measure zero (\bgref{C.22}). The paper writes $\lvert M\rvert=\dim Q$ for the parameter budget, $d(M)$
for the generic rank, $S_1=d\rho_{q_0}(T_{q_0}Q)$ for the \textbf{first-order image} and $d(M)-\dim S_1$ for the
\textbf{first-order deficit}. For definable $\rho$ on an open definable $Q$, as for every method in this paper,
\bgref{C.18}(d) shows that $d(M)=\dim\Im M$, the \textbf{expressive dimension}. The deficit vanishes exactly when $q_0$
is a point of maximal rank; the base point $\theta_0$ is then called \textbf{regular} for $M$, and an \textbf{apex} when
the deficit is positive (\thmref{Definition}{I.5}).
\end{bgentry}

\thmref{Proposition}{II.3} gives LoRA's rank at every point,
$m\operatorname{rk}A+n\operatorname{rk}B-\operatorname{rk}A\operatorname{rk}B$. At a random pair it is $r(m+n-r)=d$; at
the zero-init point $(0,A_0)$ it is $mr$, a deficit of $r(n-r)$. Since the points of lower rank have measure zero, a
float64 Jacobian rank at one random point computes $d(M)$ with probability one, up to the tolerance that decides
numerical rank. The plug workbench of the companion site (\S\ref{sec:cut-plug}) and the check script
(Appendix~\ref{app:repro}) do exactly this.

\fbAusedin{the notation (Table~\ref{tab:notation}), \thmref{Definition}{I.5}, \thmref{Proposition}{II.3},
  \thmref{Proposition}{II.12}, the plug workbench (\S\ref{sec:cut-plug}), reproducing the numbers
  (Appendix~\ref{app:repro}).}

\begin{bgentry}{C.4}{immersions, submersions and embedded submanifolds}
A smooth $\rho:Q\to\Theta$ is an \textbf{immersion} if every $d\rho_q$ is injective (rank $\dim Q$), and a
\textbf{submersion} if every $d\rho_q$ is surjective (rank $\dim\Theta$). A subset $X\subseteq\R^N$ is an
\textbf{embedded submanifold} of dimension $k$ if near each of its points it is the zero set of $N-k$ smooth functions
with linearly independent gradients. A submersion is an open map, so its image contains an open set. An
\textbf{embedding}, an injective immersion that is a homeomorphism onto its image, has an embedded submanifold as image;
an injective immersion such as the figure eight need not.
\end{bgentry}

The one-sided method $\mathrm{FA}(A_0):B\mapsto W_0+BA_0$ is an injective immersion with affine image $W_0+\{BA_0\}$
when $A_0$ has full row rank. On pairs of full rank, $(B,A)\mapsto BA$ is a submersion onto $\MM_r$ (\bgref{C.12}).
\thmref{Observation}{I.4}(b) uses openness, in the argument recalled in \bgref{A.13}. \thmref{Proposition}{IV.3} asks
for injective immersions $(\R^k,0)\to(U,\theta)$ tangent to a distribution.

\fbAusedin{\thmref{Observation}{I.4}, \thmref{Proposition}{IV.3}, integrability (\S\ref{sec:lens-integrability}).
  \textsf{Reading:} \citet[Ch.~4 and~5]{xlee2013introduction}.}

\begin{bgfact}{C.5}{constant-rank theorem}
Let $\rho:Q\to\Theta$ be smooth with $\rk d\rho_q=k$ for all $q$ near $p$. Then there are charts centred at $p$ and at
$\rho(p)$ in which $\rho$ reads $(x_1,\dots,x_{\dim Q})\mapsto(x_1,\dots,x_k,0,\dots,0)$. Hence some neighbourhood $V$
of $p$ has $\rho(V)$ an embedded $k$-dimensional submanifold, and $\rho^{-1}(\rho(p))\cap V$ an embedded submanifold of
dimension $\dim Q-k$. By \bgref{C.3} the hypothesis holds wherever the rank is maximal.
\end{bgfact}

\paragraph{The idea of the proof}
Choose coordinates in which the top-left $k\times k$ block of the Jacobian is invertible. The inverse function theorem
straightens the first $k$ output coordinates. Constant rank then forces the other outputs to depend only on those $k$,
and a second change of coordinates in the target removes them.

The paper uses this theorem without naming it. Near a point of maximal rank, $\rho$ is locally a projection onto a sheet
of dimension $d(M)$, with fibres of dimension $\lvert M\rvert-d(M)$, the gauge waste. For LoRA at a full-rank pair the
sheet is a piece of $W_0+\MM_r$ and the fibre a piece of a $\GL_r$-orbit. \thmref{Proposition}{IV.3} makes the analogous
constant-rank assumption for distributions, where Frobenius' theorem (\bgref{D.25}) takes the place of this one.

\fbAusedin{\thmref{Definition}{I.5}, \thmref{Proposition}{II.12}, \bgref{C.18}. \textsf{Reading:} \citet[Ch.~4, the rank
  theorem]{xlee2013introduction}.}

\begin{bgentry}{C.6}{fibres and fibre dimension}
The \textbf{fibre} of $\rho$ over $\theta$ is $\rho^{-1}(\theta)$, the set of parameters that give the weight $\theta$.
If $\rk d\rho_q=d(M)$, \bgref{C.5} makes the fibre through $q$ a submanifold of dimension $\lvert M\rvert-d(M)$ near
$q$. For real-analytic $\rho$ on connected $Q$ this holds for all $q$ outside a closed set of measure zero, and these
are the \textbf{generic fibres}. Fibres through points of lower rank can have a different dimension. The fibres of a
fibre bundle all look alike locally, so a map whose fibre dimension jumps is no fibre bundle.
\end{bgentry}

For LoRA on the full-rank stratum, the fibre through $(B,A)$ is the orbit $\{(Bg^{-1},gA):g\in\GL_r\}$, of dimension
$r^2=r(m+n)-r(m+n-r)$ (\bgref{D.21}). For zero-init LoRA the fibre over $\theta_0$ is $\{(B,A):BA=0\}$. It contains
$\{B=0\}$, of dimension $rn>r^2$, and is itself a union of strata, such as $\{(0,A):\rk A=r\}$
(\thmref{Corollary}{III.13}).

\fbAusedin{Exposé~\ref{sec:fibres}, \thmref{Definition}{I.5}, \thmref{Proposition}{II.12}, \thmref{Corollary}{III.13},
  Rosetta entries 6 and 11 (Exposé~\ref{sec:rosetta}), the seven borrowed words (Exposé~\ref{sec:analogy}).}

\begin{bgentry}{C.7}{fibre products, transversality and clean intersection}
For $\rho_M:Q_M\to\Theta$ and $\rho_N:Q_N\to\Theta$, the set-level \textbf{fibre product} is
$P=\{(q,q'):\rho_M(q)=\rho_N(q')\}\subseteq Q_M\times Q_N$. The maps are \textbf{transverse}, $\rho_M\pitchfork\rho_N$,
if $d\rho_M(T_qQ_M)+d\rho_N(T_{q'}Q_N)=T_\theta\Theta$ at every $(q,q')\in P$; then $P$ is an embedded submanifold of
dimension $\lvert M\rvert+\lvert N\rvert-\dim\Theta$, and it is the pullback of $\rho_M$ and $\rho_N$ in $\Diff$
(\bgref{A.11}). The intersection is \textbf{clean} if $P$ is an embedded submanifold with
$T_{(q,q')}P=\{(v,w):d\rho_Mv=d\rho_Nw\}$ at every point. Transverse maps meet cleanly, and so do affine maps.
\end{bgentry}

Transversality at $(q_0,q_0')$ needs $\lvert M\rvert+\lvert N\rvert\ge\dim\Theta$, which fails for any two small
methods, so the products the paper builds come from clean cases. The affine methods $\mathrm{FA}(A_0)$ and
$\mathrm{FB}(B_0)$, with $B_0$ of full column rank and $A_0$ of full row rank, have the linear fibre product
$\{(B_0R,RA_0)\}$ of dimension $r^2$, which is LoRA-XS \citep{baazy2024lora} (\thmref{Proposition}{I.8}). On
$\Theta=\R$, $\mathrm{LoRA}_1$ and the frozen model have the cross $\{BA=0\}$ as fibre product. It is no manifold at the
origin, and no product exists (\bgref{A.13}).

\fbAusedin{\thmref{Observation}{I.4}, \thmref{Proposition}{I.8}, \thmref{Theorem}{II.7}(d). \textsf{Reading:}
  \citet[Ch.~1]{xguillemin1974differential}; \citet[Ch.~6]{xlee2013introduction}.}

% ----------------------------------------------------------------------------------------------------------------------
\subsubsection*{Germs, cones and singular points}

\begin{bgentry}{C.8}{germs}
Two subsets $X,Y\subseteq\Theta$ have the same \textbf{germ} at $\theta_0$ if $X\cap U=Y\cap U$ for some neighbourhood
$U$ of $\theta_0$; two maps have the same germ at $q_0$ if they agree near $q_0$. The germs $(X,\theta_0)$ and $(Y,0)$
are diffeomorphic if a diffeomorphism $\varphi$ between neighbourhoods $U\ni\theta_0$ and $V\ni0$ has
$\varphi(\theta_0)=0$ and $\varphi(X\cap U)=Y\cap V$. Local properties, such as the dimension at $\theta_0$, smoothness
at $\theta_0$ and the tangent cone (\bgref{C.11}), depend only on the germ.
\end{bgentry}

The image germ of a method is the germ of $\Im M=\rho(Q)$ at $\theta_0$. It uses all of $Q$, so it can be larger than
the image of a small neighbourhood of $q_0$. For zero-init LoRA, $sBA$ with $B$ small and $A$ arbitrary is a small
update with any row space, so the germ is that of the whole cone $\theta_0+\Mr$; parameters near $(0,A_0)$ only give
updates with row space inside $\operatorname{row}A$ for some $A$ near $A_0$. For injective $W_0$ and even $r$ with $2\le
r\le n-2$, HRA's image germ at $W_0$ is the image of the germ of $\Sigma_r$ at $0$ under the embedding $\Omega\mapsto
W_0\,C(\Omega)$ of $\so(n)$ into $\R^{m\times n}$ (\bgref{C.13}, \bgref{D.7}).

\fbAusedin{the image germ (Exposé~\ref{sec:image}), the three invariants (Exposé~\ref{sec:thesis}),
  \thmref{Theorem}{II.7}, the seven axes (\S\ref{sec:atlas-axes}).}

\begin{bgentry}{C.9}{smooth and singular points}
A point $x$ of a set $X\subseteq\R^N$ is a \textbf{smooth point} if $X\cap U$ is a $C^1$ embedded submanifold for some
neighbourhood $U$ of $x$, and a \textbf{singular point} otherwise. The singular points form the \textbf{singular locus}
$\operatorname{Sing}X$. Algebraic geometry defines $\operatorname{Sing}$ through the Jacobian of defining polynomials,
and for $\Mr$ the two notions agree. Smoothness is a property of a set. The paper's regular base point and apex, defined
by a zero or positive first-order deficit (\bgref{C.3}, \thmref{Definition}{I.5}), are properties of a method at $q_0$.
\end{bgentry}

The cross $\{BA=0\}\subseteq\R^2$ is singular at the origin and smooth elsewhere. For $1\le r<\min(m,n)$,
$\operatorname{Sing}\Mr=\MM_{\le r-1}$. A split initialisation has image $(\theta_0-P)+\Mr$ with $\rk P=r$, so
$\theta_0$ is a smooth point; zero-init LoRA starts at the singular vertex. A figure-eight curve passing through its
crossing with nonzero velocity has deficit $0$ there, though its image is singular, so smooth and regular are different
notions.

\fbAusedin{\thmref{Definition}{I.5}, \thmref{Theorem}{II.2}, Figure~\ref{fig:apex}, \thmref{Theorem}{II.7}, Erlangen
  (\S\ref{sec:erlangen}). \textsf{Reading:} \citet[Lecture~14]{xharris1992algebraic}.}

\begin{bgentry}{C.10}{cones and vertices}
A set $X$ in a vector space is a \textbf{cone} with \textbf{vertex} $v$ if $v\in X$ and $v+t(x-v)\in X$ for all $x\in X$
and $t\ge0$, and \textbf{linear} if it is an affine subspace. If $\rho=\theta_0+\delta$ with
$\delta(tu,w)=t\,\delta(u,w)$ for $t\ge0$ on a block $u$ of parameters, and $Q$ is closed under $u\mapsto tu$, then $\Im
M$ is a cone with vertex $\theta_0$. A cone that is a $C^1$ embedded submanifold near its vertex $v$ equals $v+T_vX$, so
the vertex of a non-linear cone is a singular point.
\end{bgentry}

\paragraph{Why a smooth cone is linear}
Each ray from $v$ lies in $X$, so $X\subseteq v+T_vX$. Near $v$, $X$ is a submanifold of full dimension in this affine
space, so it contains a neighbourhood of $v$ there, and scaling spreads it over all of $v+T_vX$.

$\Mr$ is a closed cone with vertex $0$, and it is not linear for $1\le r<\min(m,n)$, since a rank-$r$ and a rank-one
matrix can sum to rank $r+1$. Zero-init LoRA has $\delta=sBA$, linear in $B$, so its image $W_0+\Mr$ has singular vertex
$W_0$. Its deficit $r(n-r)$ makes $\theta_0$ an apex too; the vertex is a property of the set, the apex one of the
method. Multihomogeneous updates such as LoHa's \citep{hyeonwoo2021fedpara} give cones in the same way (\bgref{D.19}).

\fbAusedin{\thmref{Definition}{I.5}, the image germ (Exposé~\ref{sec:image}), \thmref{Theorem}{II.1},
  \thmref{Theorem}{II.2}, Figure~\ref{fig:apex}, \thmref{Theorem}{II.7}.}

\begin{bgfact}{C.11}{first-order image versus image}
The \textbf{tangent cone} of a set $X$ at $x$ is the set of limits of $t_k(x_k-x)$ with $x_k\in X$, $x_k\to x$ and
$t_k>0$; for a closed cone with vertex $v$ it is $X-v$. For a method $M$: (a) $S_1$ lies in the tangent cone of $\Im M$
at $\theta_0$; (b) $\dim S_1\le d(M)$; (c) if $\Im M$ is a $C^1$ manifold of dimension $d(M)$ near $\theta_0$, then
$S_1\subseteq T_{\theta_0}\Im M$, with equality exactly when the deficit is $0$.
\end{bgfact}

For (a), $d\rho_{q_0}(\xi)$ is the velocity of the curve $\rho(q_0+t\xi)$ inside the image; for (c), $\rho$ maps a
neighbourhood of $q_0$ into the manifold. For $\mathrm T_A$-pointed LoRA ($B_0=0$, \bgref{C.19}) the tangent cone at
$\theta_0$ is $\Mr$ itself, and inside it $S_1=\{XA_0\}$, the matrices with row space in $\operatorname{row}A_0$, is a
linear space of dimension $mr<r(m+n-r)$. Figure~\ref{fig:apex}(a) draws $S_1$ as one ruling of the cone, a line through
the apex. At a split pointing case (c) applies, and $S_1=T_P\MM_r$, the tangent plane of Figure~\ref{fig:apex}(b).

\fbAusedin{\thmref{Definition}{I.5}, \thmref{Observation}{I.6}, \thmref{Theorem}{II.2}, Figure~\ref{fig:apex}.
  \textsf{Reading:} \citet[Lecture~20]{xharris1992algebraic}.}

% ----------------------------------------------------------------------------------------------------------------------
\subsubsection*{Matrix varieties}

\begin{bgentry}{C.12}{the determinantal variety}
$\Mr=\{X\in\R^{m\times n}:\rk X\le r\}$, the \textbf{determinantal variety}, is the zero set of all $(r+1)\times(r+1)$
minors, a closed cone with vertex $0$. Let $1\le r<\min(m,n)$. The matrices of rank exactly $r$ form an open dense
subset $\MM_r$ of $\Mr$, an embedded submanifold of dimension $r(m+n-r)$, and $\operatorname{Sing}\Mr=\MM_{\le r-1}$. If
$X\in\MM_r$ has column space $U$ and row space $V$, and $X=BA$ with $B\in\R^{m\times r}$ and $A\in\R^{r\times n}$, then
\[T_X\MM_r=\{\dot BA+B\dot A\}=\{D:P_{U^\perp}DP_{V^\perp}=0\}.\]
Also $\Mr-\Mr=\MM_{\le\min(2r,m,n)}$.
\end{bgentry}

Zero-init LoRA has image $W_0+\Mr$, of dimension $130{,}816$ for a $4096\times4096$ matrix at $r=16$. On full-rank
pairs, $(B,A)\mapsto BA$ is a submersion onto $\MM_r$ whose fibres are $\GL_r$-orbits, which identifies $\MM_r$ with the
quotient $(\R^{m\times r}_*\times\R^{r\times n}_*)/\GL_r$ (\bgref{D.22}). The difference formula is why no pointing of
$\mathrm{LoRA}_r$ reaches an update of rank above $2r$ (\bgref{B.14}, \thmref{Theorem}{II.1}).

\fbAusedin{the notation (Table~\ref{tab:notation}), translates of one shape (\S\ref{sec:strata-translates}),
  \thmref{Theorem}{II.1}, \thmref{Theorem}{II.2}, \thmref{Proposition}{II.8}, Exposé~II.C (\S\ref{sec:cut}),
  \thmref{Proposition}{III.4}. \textsf{Reading:} \citet[Lecture~9]{xharris1992algebraic}.}

\begin{bgentry}{C.13}{the skew-determinantal cone}
Every skew-symmetric matrix has even rank. For even $r$ with $2\le r\le n-2$, the \textbf{skew-determinantal cone}
$\Sigma_r=\{\Omega\in\so(n):\rk\Omega\le r\}$ is a closed cone with vertex $0$ and dimension
$r(n-r)+r(r-1)/2=rn-r(r+1)/2$, counting a range in $\Gr(r,n)$ and a nondegenerate skew form on it. Its smooth points are
the matrices of rank exactly $r$, and $\operatorname{Sing}\Sigma_r=\Sigma_{r-2}$. It is not linear, since skew matrices
of rank $r$ and $2$ can sum to rank $r+2\le n$, so its vertex is singular.
\end{bgentry}

For injective $W_0$ and even $r$, HRA's image near $W_0$ consists of the $W_0R$ with $R\in\SO(n)$ near $I$ and
$\rk(R-I)\le r$. The Cayley map satisfies $\rk(C(\Omega)-I)=\rk\Omega$ (\bgref{D.7}), so the image germ at $W_0$ is the
image of the germ of $\Sigma_r$ at $0$ under the embedding $\Omega\mapsto W_0\,C(\Omega)$, $W_0$ is a singular vertex,
and $d(\mathrm{HRA}_r)=\dim\Sigma_r$. At the paired initialisation $u_{2i-1}=u_{2i}=w_i$ ($i\le k=r/2$, with the $w_i$
independent), put $U=\spanop\{w_1,\dots,w_k\}$, so $\dim U=k$. Then $S_1=W_0\cdot\mathcal S$ with
\[\mathcal S=\{\Omega\in\so(n):P_{U^\perp}\Omega P_{U^\perp}=0\}=U\wedge\R^n\]
(\bgref{D.5}), a linear space of skew matrices of rank at most $2k=r$ and of dimension $kn-k(k+1)/2<d(\mathrm{HRA}_r)$.
It lies inside $\Sigma_r$, so $S_1$ lies inside the tangent cone $W_0\cdot\Sigma_r$ (\thmref{Theorem}{II.7}(b)).

\fbAusedin{\thmref{Theorem}{II.7}, HRA is a meet (\S\ref{sec:erlangen-hra}).}

\begin{bgentry}{C.14}{Segre and toric cones}
The \textbf{Segre cone} in $\R^p\otimes\R^q\cong\R^{p\times q}$ is the set $\{a\otimes b\}=\MM_{\le1}$ of rank-one
tensors, a closed cone of dimension $p+q-1$ that spans the whole space and is non-linear when $p,q\ge2$. It is the
simplest \textbf{toric cone}, the closure of the image of a torus $(\R^\times)^k$ under a monomial map, here
$(a,b)\mapsto(a_ib_j)_{ij}$. Since rank-one matrices span $\R^{p\times q}$, $k$ generic rank-one matrices span
$\min(k,pq)$ dimensions.
\end{bgentry}

Under the Van Loan--Pitsianis rearrangement (\bgref{B.11}), KronA's \citep{edalati2022krona} update $A\otimes B$ becomes
$\vect(A)\,\vect(B)\T$, so its image is a linear copy of a Segre cone (\thmref{Proposition}{II.12}). VeRA's
\citep{kopiczko2023vera} update $\Lambda_bB\Lambda_dA$ has entries $\sum_k(b_id_k)B_{ik}A_{kj}$, a fixed linear function
of the rank-one matrix $bd\T$. This function is injective when $B$ has no zero entry and $A$ has full row rank, so the
image is a toric cone of dimension $m+r-1$ with vertex $\theta_0$, as VeRA's entry in Exposé~\ref{sec:examples} records.
\thmref{Proposition}{II.8} uses the spanning fact for DoRA \citep{liu2024dora}, with $m$ points of the Segre cone in
$\R^{m-r}\otimes\R^{n-r}$.

\fbAusedin{\thmref{Proposition}{II.8}, \thmref{Proposition}{II.12}, \hyperref[ex:vera]{VeRA}
  (Exposé~\ref{sec:examples}). \textsf{Reading:} \citet[Lecture~2]{xharris1992algebraic}; \citet[Ch.~1]{xcox2011toric}.}

\begin{bgentry}{C.15}{the Grassmannian}
The \textbf{Grassmannian} $\Gr(r,n)$ is the set of $r$-dimensional linear subspaces of $\R^n$, a compact manifold of
dimension $r(n-r)$. One model is the full-row-rank $r\times n$ matrices up to $A\mapsto gA$, $g\in\GL_r$, through
$A\mapsto\operatorname{row}A$; another is the orthogonal projectors $P=P\T=P^2$ with $\tr P=r$. Near $V$, the graphs of
linear maps $V\to V^\perp$ give a chart onto $\operatorname{Hom}(V,V^\perp)\cong\R^{r(n-r)}$. $\Gr(r,n)$ carries a
unique $\Ort(n)$-invariant probability measure, the uniform or Haar measure, and the row space of a Gaussian $A$ has
this law (\bgref{D.14}).
\end{bgentry}

$\mathrm T_A$-pointed LoRA is classified up to isomorphism by $\operatorname{row}A_0\in\Gr(r,n)$
(\thmref{Theorem}{II.2}), and its deficit $r(n-r)$ equals $\dim\Gr(r,n)$. First steps keep the row space of the update
inside $\operatorname{row}A_0$, and the missing directions are those that tilt it. For injective $W_0$ and even $r$ with
$2\le r\le n-2$, HRA's image is the union of the orbits $W_0\cdot\SO(V)$ over $V\in\Gr(r,n)$, and its dimension
$r(n-r)+r(r-1)/2$ counts a subspace and a rotation inside it. LoRA-FA \citep{zhang2023lora}, which trains $B$ with $A_0$
frozen, sees $A_0$ only through $V=\operatorname{row}A_0$ once its gradient is corrected as in
\thmref{Proposition}{III.4}(c).

\fbAusedin{\thmref{Theorem}{II.2}, the checks (Appendix~\ref{app:repro}), \thmref{Theorem}{II.7}, LOFT
  (\S\ref{sec:erlangen-loft}) and the LOFT example there, \thmref{Proposition}{III.4}, \thmref{Corollary}{III.13}.
  \textsf{Reading:} \citet[Ch.~1]{xlee2013introduction}; \citet[Lecture~6]{xharris1992algebraic}.}

% ----------------------------------------------------------------------------------------------------------------------
\subsubsection*{Tame sets and genericity}

\begin{bgfact}{C.16}{images of smooth maps need not be manifolds}
Even a polynomial map can have an image that is no manifold. (a) \emph{Cusps.} $t\mapsto(t^2,t^3)$ is injective with
image $\{y^2=x^3\}$, which is not $C^1$ at $0$. (b) \emph{Crossings and cones.} The cross $\{BA=0\}\subseteq\R^2$ and
the cone $\Mr$ are singular at the origin. (c) \emph{Missing limits.} $(x,y)\mapsto(x,xy)$ has image
$\{u\ne0\}\cup\{(0,0)\}$ in the $(u,v)$-plane, neither open nor closed. (d) \emph{Dense images.} The line of irrational
slope is an injective immersion into the torus with dense image.
\end{bgfact}

PEFT images show (b) routinely, as in the LoRA cone and HRA's skew cone, so the paper measures them with tools that need
no manifold structure, such as germs (\bgref{C.8}), tangent cones (\bgref{C.11}) and the dimension theory of
semialgebraic and definable sets (\bgref{C.17}, \bgref{C.18}). Whether LoHa's Hadamard image is closed, the question
raised by (c), is open.

\fbAusedin{\thmref{Observation}{I.4}, \thmref{Theorem}{II.2}, \thmref{Theorem}{II.7}, the open problems
  (Exposé~\ref{sec:analogy}).}

\begin{bgentry}{C.17}{semialgebraic and definable sets}
A \textbf{semialgebraic set} is a finite union of sets $\{x\in\R^N:p(x)=0,\ q_1(x)>0,\dots,q_k(x)>0\}$ with polynomials
$p,q_i$, and a map is semialgebraic if its graph is. A set is \textbf{definable} in $\R_{\exp}$ if a first-order formula
defines it from polynomials with real coefficients, $\exp$, equality, order, Boolean operations and quantifiers over
$\R$. \citet{xwilkie1996model} proved that $\R_{\exp}$ is \textbf{o-minimal}, meaning every definable subset of $\R$ is
a finite union of points and intervals. Semialgebraic sets are the definable sets that need no $\exp$.
\end{bgentry}

LoRA, DoRA (whose row norms involve square roots), HRA and exact-Cayley OFT \citep{qiu2023controlling} are built from
polynomials, quotients and square roots, so their images are semialgebraic by \bgref{C.18}(a). A softmax brings in
$\exp$, and maps that contain one are definable in $\R_{\exp}$.

\fbAusedin{where the category theory works (Exposé~\ref{sec:thesis}), the arrow rules (\S\ref{sec:atlas-arrows}).
  \textsf{Reading:} \citet[Ch.~2]{xbochnak1998real}; \citet{xwilkie1996model}.}

\begin{bgfact}{C.18}{Tarski--Seidenberg and definable dimension}
(a) The image of a semialgebraic set under a semialgebraic map is semialgebraic (Tarski--Seidenberg); images of
definable sets under definable maps are definable. (b) Every definable set $X$ is a finite disjoint union of $C^1$
embedded submanifolds, and $\dim X$, the largest of their dimensions, does not depend on the choice. (c) $X\subseteq Y$
implies $\dim X\le\dim Y$, $\dim f(X)\le\dim X$ for definable $f$, and $\dim(\overline X\setminus X)<\dim X$ for
nonempty $X$. (d) For a definable $C^1$ map $\rho$ on an open definable $Q\subseteq\R^k$, $\dim\rho(Q)=\max_q\rk
d\rho_q$.
\end{bgfact}

\paragraph{The idea of (d)}
Near a point of maximal rank, \bgref{C.5} gives a sheet of that dimension inside $\rho(Q)$. Conversely, a definable
section of $\rho$, $C^1$ on a top-dimensional piece of $\rho(Q)$, composes with $\rho$ to the identity there, so the
rank is at least that dimension.

Part (d) is the identity $d(M)=\dim\Im M$ of \thmref{Definition}{I.5}. Part (c) is the monotonicity behind test O2
(\S\ref{sec:atlas-tests}), since an arrow $M\to N$ gives $\Im M\subseteq\Im N$ and hence $d(M)\le d(N)$; for instance
$\mathrm{LoRA}_r\not\to\text{LoRA-FA}$ because $r(m+n-r)>mr$.

\fbAusedin{\thmref{Definition}{I.5}, the arrow rules (\S\ref{sec:atlas-arrows}), where the category theory works
  (Exposé~\ref{sec:thesis}). \textsf{Reading:} \citet[Ch.~3 and~4]{xvandendries1998tame};
  \citet[Ch.~2]{xbochnak1998real}.}

\begin{bgentry}{C.19}{stratification and local dimension}
A \textbf{stratification} of a set $X$ is a partition into finitely many embedded submanifolds, the \textbf{strata},
such that the closure of each stratum is a union of strata. The \textbf{local dimension} $\dim_xX$ is the dimension of
the germ of $X$ at $x$, the largest dimension of a stratum whose closure contains $x$. $X$ is \textbf{pure-dimensional}
if $\dim_xX$ is the same at every point. Definable sets admit stratifications, so these notions make sense for every
image in this paper.
\end{bgentry}

The cross $\{ba=0\}\subseteq\R^2$ is stratified as the origin and the four open half-axes, and it is pure of dimension
$1$. Likewise $\Mr=\MM_0\sqcup\MM_1\sqcup\dots\sqcup\MM_r$ with $\overline{\MM_k}=\MM_{\le k}$, and it is pure of
dimension $r(m+n-r)$. For $1\le r<\min(m,n)$, \thmref{Theorem}{II.2} singles out three strata of full-rank LoRA
pointings, $\mathrm T_A$ ($B_0=0$, $\rk A_0=r$), $\mathrm T_B$ ($A_0=0$, $\rk B_0=r$) and the open dense $\mathrm S$
($\rk B_0A_0=r$), and the fibre $\{BA=0\}$ over $\theta_0$ is itself stratified (\thmref{Corollary}{III.13}). The HRA
fibre product of \thmref{Theorem}{II.7}(d) has local dimension $d(\mathrm{HRA}_r)+r^2=rn+r(r-1)/2$ over generic points
of the meet and contains $\{(B,0,I)\}$, of dimension $mr$, over $W_0$, so it is not pure-dimensional when
$mr>rn+r(r-1)/2$.

\fbAusedin{\thmref{Theorem}{II.2}, the three strata (\S\ref{sec:strata-three}), \thmref{Theorem}{II.7},
  \thmref{Corollary}{III.13}, the open problems (Exposé~\ref{sec:analogy}), the seven borrowed words
  (Exposé~\ref{sec:analogy}). \textsf{Reading:} \citet[Ch.~4]{xvandendries1998tame}; \citet[Ch.~9]{xbochnak1998real}.}

\begin{bgentry}{C.20}{moduli sets and moduli spaces}
Classifying objects up to isomorphism gives a set of classes. A \textbf{moduli space} also carries a topology or
geometry for which continuous or smooth families of objects correspond to continuous or smooth maps into it, ideally
with a universal family over it from which every family is pulled back. A disjoint union of spaces $X\sqcup Y$ has the
topology in which $X$ and $Y$ are both open and closed.
\end{bgentry}

\thmref{Theorem}{II.2} classifies the three strata of LoRA pointings, as a set, by $\Gr(r,n)\sqcup\Gr(r,m)\sqcup\MM_r$.
The quotient topology does not split this way, since $(\varepsilon B',A_0)\in\mathrm S$ tends to $(0,A_0)\in\mathrm T_A$
as $\varepsilon\to0$, so $\mathrm T_A$ lies in the closure of $\mathrm S$. No universal family is built, and the paper
uses ``moduli'' as an analogy.

\fbAusedin{\thmref{Theorem}{II.2}, the Rosetta stone (Exposé~\ref{sec:rosetta}), the seven borrowed words
  (Exposé~\ref{sec:analogy}). \textsf{Reading:} \citet[Lecture~4]{xharris1992algebraic}.}

\begin{bgentry}{C.21}{generic properties}
A property of points of $\R^N$ holds \textbf{generically} if the set where it fails lies in $\{f=0\}$ for some nonzero
polynomial $f$. The set $\{f=0\}$ is closed, nowhere dense and of Lebesgue measure zero, so a generic property holds
with probability one at a random point whose law has a density, such as a Gaussian or uniform initialisation. For points
of a parametrised set, such as a Segre cone, ``generic'' refers to generic parameters. To prove genericity it suffices
to write failure as the vanishing of polynomials and exhibit one point where one of them is nonzero. A weaker notion,
\textbf{almost every} point, asks only that the failure set have Lebesgue measure zero, as a countable union of proper
algebraic subsets does.
\end{bgentry}

\thmref{Theorem}{II.2}(d) uses a generic rank-one update, \thmref{Proposition}{II.8} gives DoRA's dimension
$\min\big(mn,\,r(m+n-r)+m\big)$ for generic $W_0$, and \thmref{Proposition}{V.2} notes that QLoRA's
\citep{dettmers2023qlora} defect $e=\kappa W-W$ has full rank for almost every $W$, in the weaker sense. The quantiser
$\kappa$ is defined cell by cell rather than by one polynomial, so the argument runs cell by cell, and the failure set
is a countable union of proper algebraic subsets of quantisation cells, which has measure zero. A pretrained $W_0$ is
not random, and structured weights can lie in the exceptional set. For a rank-one $W_0=ac\T$ with all $a_i\ne0$ and
$r<n$, \thmref{Proposition}{II.8} finds DoRA's dimension $r(m+n-r)+m-r$, strictly below the generic value exactly when
$r<m$ and $r\le n-2$.

\fbAusedin{\thmref{Theorem}{II.2}, \thmref{Proposition}{II.8}, \thmref{Theorem}{II.11}, \thmref{Proposition}{II.12},
  \thmref{Proposition}{V.2}.}

\begin{bgentry}{C.22}{real-analytic functions and their zero sets}
A function is \textbf{real-analytic} on an open set if near every point it equals a convergent power series.
Polynomials, $\exp$, softmax, $\tanh$, GELU and SiLU are real-analytic, and so are sums, products, compositions and
quotients with nonvanishing denominator. On a connected open set, a real-analytic function that vanishes on a nonempty
open subset vanishes identically, and the zero set of one that is not identically zero is closed, nowhere dense and of
measure zero.
\end{bgentry}

\thmref{Lemma}{VI.2} compares Taylor coefficients at $0$. A nonzero scale $c$ with $\sigma(ct)=c\,\sigma(t)$ must
satisfy $ca_k=c^ka_k$ for every coefficient $a_k$ of $\sigma$, so $c=1$ for GELU and SiLU, whose $a_2\ne0$.
\thmref{Theorem}{VI.5}(f) uses the zero-set fact. A prefixed head is analytic in the prefix, and one prefix with a
nonzero second derivative makes the prefixes with affine output a nowhere-dense set. The same fact makes the
maximal-rank set of an analytic $\rho$ open and dense (\bgref{C.3}).

\fbAusedin{\thmref{Lemma}{VI.2}, \thmref{Theorem}{VI.5}, Rosetta entry 5 (Exposé~\ref{sec:rosetta}). \textsf{Reading:}
  \citet{xkrantz2002primer}.}

% ======================================================================================================================
\subsection*{Part D. Groups, symmetry and conservation}
\phantomsection\label{app:bg-d}\addcontentsline{toc}{subsection}{Part D. Groups, symmetry and conservation}

Part D turns to symmetry. Most PEFT parametrisations have symmetries, since different trainable parameters can give the
same weights and some transformations of weight space carry a method's image to itself. This part collects the group
theory the paper uses to describe those symmetries, the invariants they leave untouched, and the quantities they
conserve under gradient flow.

% ----------------------------------------------------------------------------------------------------------------------
\subsubsection*{Groups and their Lie algebras}

\begin{bgentry}{D.1}{Lie groups; the matrix groups of this paper}
A \textbf{Lie group} is a group that is also a smooth manifold, with smooth multiplication and inversion. Every closed
subgroup of $\GL_n(\R)$ is a Lie group (Cartan's closed-subgroup theorem), and all groups below are of this kind. Its
dimension is its dimension as a manifold.

\smallskip
{\centering\small\renewcommand{\arraystretch}{1.12}
\begin{tabularx}{\linewidth}{@{}llX@{}}
  \toprule
  \headrow \textbf{group} & \textbf{dimension} & \textbf{elements}\\
  \midrule
  $\GL_r$ & $r^2$ & invertible $r\times r$ matrices\\
  $\Ort(n)$, $\SO(n)$ & $n(n-1)/2$ & $R\T R=I$; for $\SO(n)$ also $\det R=1$\\
  $\CO(n)$ & $n(n-1)/2+1$ & $cR$ with $c>0$, $R\in\Ort(n)$, so $\CO(n)=\R_{>0}\cdot\Ort(n)$; equivalently
    $g\T g\in\R_{>0}I$\\
  $\mathrm{Sp}(2r,\R)$ & $r(2r+1)$ & $M\T JM=J$, with $J=\begin{pmatrix}0&I_r\\-I_r&0\end{pmatrix}$\\
  $T_m$ & $m$ & invertible diagonal $m\times m$ matrices (\bgref{D.2})\\
  \bottomrule
\end{tabularx}\par}
\end{bgentry}

LoRA's gauge is $\GL_r$, acting on factors by $g\cdot(B,A)=(Bg^{-1},gA)$. Exactly orthogonal OFT keeps the weight inside
the orbit $W_0\cdot\Ort(n)$, and HRA with even $r$ keeps it inside $W_0\cdot\SO(n)$. $\CO(n)$ is exactly the input-side
covariance group (\bgref{D.17}) of the PiSSA \citep{meng2024pissa}, MiLoRA \citep{wang2024milora} and paper-form CorDA
\citep{yang2024corda} splits. LoCO's \citep{nguyen2026loco} skew generator satisfies $UV\T-VU\T=ZJZ\T$ with $Z=[U\ V]$,
so it is unchanged under $Z\mapsto ZM$ for every $M\in\mathrm{Sp}(2r,\R)$.

\fbAusedin{the notation (Table~\ref{tab:notation}), \thmref{Definition}{I.5}, \thmref{Theorem}{II.5},
  \thmref{Proposition}{II.16}, LOFT (\S\ref{sec:erlangen-loft}), \thmref{Theorem}{III.5}, \thmref{Theorem}{V.3}.
  \textsf{Reading:} \citet[Ch.~1]{xhall2015lie}.}

\begin{bgentry}{D.2}{tori, diagonal monoids and saturation}
The \textbf{diagonal torus} $T_m=\{\diag(d):d_i\ne0\}\cong(\R^\times)^m$ is a commutative Lie group with $2^m$ connected
components, and $T_m^+=\{\diag(d):d_i>0\}\cong\R_{>0}^m$ is the component of the identity. $\mathrm{Diag}_m$, the set of
all diagonal matrices including singular ones, is a monoid under multiplication (\bgref{D.10}) whose invertible elements
form $T_m$. Identifying $\R^2$ with $\mathbb C$, the scaled rotations $a\,\mathrm{Rot}(\theta)$ with $a>0$ form
$\CO^+(2)\cong\mathbb C^\times$, through $a\,\mathrm{Rot}(\theta)\mapsto ae^{i\theta}$. If a group or monoid $G$ acts on
a space containing a set $X$, the \textbf{saturation} $G\cdot X=\{g\cdot x:g\in G,\ x\in X\}$ is the smallest $G$-stable
set containing $X$.
\end{bgentry}

(IA)$^3$ \citep{liu2022few} on keys and values maps $\ell\mapsto\diag(\ell)W_0$, so its image $\mathrm{Diag}_m\cdot W_0$
is the closure of the torus orbit $T_m\cdot W_0$. DoRA's image is the saturation $\mathrm{Diag}_m\cdot(W_0+\Mr)$ of
LoRA's cone. RoAd$_1$ \citep{liao20243} is the left action of $\CO^+(2)^{m/2}\cong(\mathbb C^\times)^{m/2}$ on pairs of
output rows, which makes it (IA)$^3$ over $\mathbb C$. LoHa's gauge contains a two-sided diagonal torus.

\fbAusedin{\thmref{Theorem}{II.5}, DoRA (\S\ref{sec:erlangen-dora}), \thmref{Proposition}{II.8},
  \thmref{Proposition}{II.9}, \thmref{Proposition}{II.12}, the hyperparameter torus (\S\ref{sec:fibres-torus}).}

\begin{bgentry}{D.3}{permutations, monomial matrices, semidirect products, normalisers}
$S_d$ is the group of permutations of $d$ objects, realised by permutation matrices. A matrix is \textbf{monomial} if
each row and each column has exactly one nonzero entry. Every monomial matrix factors uniquely as $DP$ with $D\in T_m$
and $P$ a permutation matrix, and $T_m$ is normal in the group $\mathrm{Mon}_m$ of monomial matrices, so
$\mathrm{Mon}_m=T_m\rtimes S_m$. Here a \textbf{semidirect product} $N\rtimes H$ is a group with a normal subgroup $N$
and a subgroup $H$ such that $N\cap H=\{e\}$ and $NH$ is the whole group. The \textbf{signed permutations}
$B_r=\{\pm1\}^r\rtimes S_r=\mathrm{Mon}_r\cap\Ort(r)$ are the monomial matrices with entries $\pm1$, and $\lvert
B_r\rvert=2^rr!$. The \textbf{normaliser} of a subset $H\subseteq G$ is $N_G(H)=\{g\in G:gHg^{-1}=H\}$.
\end{bgentry}

A monomial $g$ satisfies $g\,\mathrm{Diag}_m\,g^{-1}=\mathrm{Diag}_m$, and in $\GL_m$ the normaliser of $T_m$ is exactly
$\mathrm{Mon}_m$. This puts $\mathrm{Mon}_m\times\GL_n$ inside the covariance group (\bgref{D.17}) of output scalings
such as (IA)$^3$ and DoRA's magnitude. Adam \citep{xkingma2015adam} acts coordinatewise, so on LoRA's full-rank locus with
channel-uniform hyperparameters (one learning rate, $\varepsilon$ and decay per factor) it is equivariant, of LoRA's
gauge, only under $B_r\subset\Ort(r)\subset\GL_r$ (\thmref{Theorem}{III.12}). By \thmref{Lemma}{VI.2}, an invertible linear map that
commutes with a non-affine $C^2$ activation is monomial.

\fbAusedin{DoRA (\S\ref{sec:erlangen-dora}), \thmref{Proposition}{II.8}, \thmref{Proposition}{II.9},
  \thmref{Theorem}{III.12}, \thmref{Corollary}{III.13}, \thmref{Proposition}{V.4}, \thmref{Lemma}{VI.2}, the privileged
  basis (\S\ref{sec:merge-basis}).}

\begin{bgentry}{D.4}{connectedness and the identity component}
The \textbf{identity component} $G^\circ$ of a Lie group $G$ is the connected component containing $e$. It is an open
normal subgroup. $\GL_n$ has two components, told apart by the sign of the determinant, and $\GL^+(n)=\{\det>0\}$ is the
identity component; likewise $\SO(n)\subset\Ort(n)$ and $T_m^+\subset T_m$. A subset $U$ is \textbf{symmetric} if
$U^{-1}=U$. Every neighbourhood of $e$ contains a symmetric one, and in a connected group a symmetric neighbourhood $U$
of $e$ generates the group as a monoid, so every element is a finite product $u_1\cdots u_k$ with $u_i\in U$.
\end{bgentry}

In merge-and-restart of a method of action type (\bgref{D.12}), each cycle applies one element of a set $S$. If $S$
contains a neighbourhood of the identity in a connected group $H$, repeated merging reaches all of $H$
(\thmref{Theorem}{V.3}(b)). HF's \citep{xmangrulkar2022peft} Cayley--Neumann OFT factors have $\det R\ge0$, so the
invertible ones lie in $\GL^+(b)$; they generate $\prod\GL^+(b)$ as a group, or $\prod\CO^+(2)$ when $b=2$.

\fbAusedin{complete invariants (\S\ref{sec:erlangen-invariants}), \thmref{Theorem}{V.3}, the corollaries of
  \thmref{Theorem}{V.3} (\S\ref{sec:base-rebase}).}

\begin{bgentry}{D.5}{Lie algebras; the commutator; $\gl$, $\so$, $\Sym$, $\wedge$}
The \textbf{Lie algebra} of a matrix group $G\subseteq\GL_n$ is $\mathfrak g=\{X:e^{tX}\in G\text{ for all }t\in\R\}$,
the tangent space of $G$ at $I$. It is a vector space closed under the \textbf{commutator} $[X,Y]=XY-YX$. The paper uses
$\gl_r$ (all $r\times r$ matrices, for $\GL_r$), $\so(n)$ (skew matrices $\Omega\T=-\Omega$, for $\Ort(n)$ and
$\SO(n)$), the diagonal matrices (for $T_m$), $\R I\oplus\so(n)$ (for $\CO(n)$) and $\mathfrak{sp}(2r)=\{X:X\T
J+JX=0\}$. For a subspace $V\subseteq\R^n$, $\so(V)$ consists of the skew maps of $V$ extended by $0$ on $V^\perp$. The
symmetric matrices $\Sym_r$ do not form a Lie algebra, since $[S,S']$ is skew. Write $v\wedge w=vw\T-wv\T\in\so(n)$,
$U\wedge V$ for the span of the $u\wedge v$ with $u\in U$, $v\in V$, and $L_{ab}=e_a\wedge e_b=E_{ab}-E_{ba}$. Then
$[L_{ac},L_{cb}]=L_{ab}$ for distinct $a,b,c$.
\end{bgentry}

For injective $W_0$, at HRA's paired initialisation $u_{2i-1}=u_{2i}=w_i$ ($i\le k=r/2$, with independent $w_i$ and
$U=\spanop\{w_i\}$ of dimension $k$), the first-order image is $S_1=W_0\cdot\{\Omega\in\so(n):P_{U^\perp}\Omega
P_{U^\perp}=0\}=W_0\cdot(U\wedge\R^n)$, a subspace of $W_0\cdot\so(n)$ of dimension $kn-k(k+1)/2$ (\bgref{C.13}). LOFT
\citep{zhao2026loft} trains a skew $E\in\so(r)$ inside a fixed $r$-dimensional input subspace. The bracket identity lets
rotations of overlapping coordinate blocks generate rotations of their union (\bgref{D.8}).

\fbAusedin{\thmref{Theorem}{II.7}, the LOFT example and LOFT fixes one leaf (\S\ref{sec:erlangen-loft}),
  \thmref{Theorem}{III.5}, \thmref{Theorem}{V.3}. \textsf{Reading:} \citet[Ch.~3]{xhall2015lie}.}

\begin{bgentry}{D.6}{the exponential map; one-parameter subgroups}
For $X\in\mathfrak g$, the curve $t\mapsto e^{tX}$ of matrix exponentials is a \textbf{one-parameter subgroup}, a smooth
homomorphism $(\R,+)\to G$, and every continuous one-parameter subgroup has this form for a unique $X$, its
\textbf{generator}. The exponential map $\mathfrak g\to G$ has derivative $\id$ at $0$, so it maps a neighbourhood of
$0$ diffeomorphically onto a neighbourhood of $I$. It is onto for $\SO(n)$ and not onto for $\GL^+(n)$, since
$\diag(-1,-2)$ is $e^X$ for no real $X$.
\end{bgentry}

The one-parameter subgroup $t\mapsto(Be^{-tX},e^{tX}A)$ of LoRA's gauge moves along a fibre without changing $BA$.
\thmref{Theorem}{III.5} builds its isometries as $k=e^{\tau\xi}$ from generators $\xi$ with $\Lambda^{-1}\xi$
antisymmetric. LOFT may use the exponential as its map $\so(r)\to\SO(r)$, and the Kempf--Ness argument (\bgref{D.30})
differentiates the norm along $e^{t\xi}$.

\fbAusedin{\thmref{Theorem}{III.5}, \thmref{Proposition}{III.7}, LOFT (\S\ref{sec:erlangen-loft}),
  \thmref{Theorem}{V.3}. \textsf{Reading:} \citet[Ch.~2]{xhall2015lie}.}

\begin{bgentry}{D.7}{the Cayley map and its Neumann truncation}
For skew $\Omega$ the eigenvalues are imaginary, so $I-\Omega$ is invertible, and the \textbf{Cayley map} is
\[C(\Omega)=(I+\Omega)(I-\Omega)^{-1}.\]
It lands in $\SO(n)$, and $C(\Omega)-I=2\Omega(I-\Omega)^{-1}$, so $\rk(C(\Omega)-I)=\rk\Omega$. Its expansion
$C(\Omega)=I+2\Omega+2\Omega^2+\cdots$ shows that its derivative at $0$ is $2\,\id$. It is a diffeomorphism from
$\so(n)$ onto the open dense set of rotations with no eigenvalue $-1$, with inverse $R\mapsto(R-I)(R+I)^{-1}$, so it
misses, for instance, $\diag(-1,-1,1,\dots,1)$.
\end{bgentry}

Exact-Cayley OFT and block OFT are orthogonal, so they preserve $K_{\mathrm{out}}=WW\T$ (\bgref{D.15}) and the spectrum.
HF PEFT's default (since version 0.18) truncates $(I-Q)^{-1}$ to $I+Q+Q^2+Q^3$, which gives
$R=C(Q)(I-Q^4)$. As $Q^4$ is symmetric and
commutes with $C(Q)$, $R\T R=(I-Q^4)^2\ne I$, and for small $Q$ each merge shrinks $K_{\mathrm{out}}$. The rank identity
is how \thmref{Theorem}{II.7}(c) identifies HRA's image germ with the cone of skew matrices of rank at most $r$.

\fbAusedin{the thesis (Exposé~\ref{sec:thesis}), \thmref{Corollary}{II.6}, \thmref{Theorem}{II.7}, the Gram lab of the
  companion site (\S\ref{sec:erlangen-invariants}), the corollaries of \thmref{Theorem}{V.3} (\S\ref{sec:base-rebase}).}

\begin{bgentry}{D.8}{Lie subgroups and subalgebras; the generated Lie algebra}
A \textbf{Lie subgroup} of $G$ is a subgroup with a Lie group structure for which the inclusion is an injective
immersion; it need not be closed. Connected Lie subgroups $H\subseteq G$ correspond one to one to Lie subalgebras
$\mathfrak h\subseteq\mathfrak g$, with $H$ the subgroup generated by $e^{\mathfrak h}$. The Lie algebra
\textbf{generated} by subalgebras $\mathfrak h_1,\dots,\mathfrak h_p$, written $\operatorname{Lie}\langle\mathfrak
h_1,\dots,\mathfrak h_p\rangle$, is the smallest subalgebra containing them, the span of their elements and all iterated
brackets of them.
\end{bgentry}

Block OFT, BOFT \citep{liu2023parameter} factors and Givens rotations \citep{ma2024parameter} use $\so(I_j)$ for
coordinate blocks $I_j\subseteq\{1,\dots,n\}$. Since $[L_{ac},L_{cb}]=L_{ab}$, two overlapping blocks generate $\so$ of
their union, and $\operatorname{Lie}\langle\so(I_j)\rangle_j=\bigoplus_C\so(C)$, summed over the connected components
$C$ of the hypergraph with vertices $\{1,\dots,n\}$ and hyperedges $I_j$. Two vertices share a component when a chain of
hyperedges, each meeting the next, joins them. Butterfly factors connect this hypergraph only at full depth.

\fbAusedin{\thmref{Theorem}{V.3}, Figure~\ref{fig:rebase}(b), the corollaries of \thmref{Theorem}{V.3}
  (\S\ref{sec:base-rebase}), LOFT fixes one leaf (\S\ref{sec:erlangen-loft}). \textsf{Reading:}
  \citet[Ch.~19]{xlee2013introduction}; \citet[Ch.~5]{xhall2015lie}.}

\begin{bgfact}{D.9}{Yamabe's theorem}
Every path-connected subgroup of a Lie group is a Lie subgroup. Consequently, if $H_1,\dots,H_p$ are connected Lie
subgroups of $G$, the subgroup they generate is the connected Lie subgroup $H$ with
$\operatorname{Lie}(H)=\operatorname{Lie}\langle\mathfrak h_1,\dots,\mathfrak h_p\rangle$ \citep{xyamabe1950arcwise}; a
short proof is in \citet{xgoto1969arcwise}.
\end{bgfact}

The subgroup generated by the $H_i$ is path-connected, because products of paths are paths, so the theorem makes it a
Lie subgroup and the correspondence of \bgref{D.8} names it. With \bgref{D.4}, this is why a re-merged action method
whose cycles contain a neighbourhood of the identity in each $H_i$ reaches exactly the connected subgroup with Lie
algebra $\operatorname{Lie}\langle\mathfrak h_1,\dots,\mathfrak h_p\rangle$, and why exact-Cayley block-orthogonal
factors reach $\prod_C\SO(C)$.

\fbAusedin{\thmref{Theorem}{V.3}, the corollaries of \thmref{Theorem}{V.3} (\S\ref{sec:base-rebase}).}

\begin{bgentry}{D.10}{monoids, units, monoid orbits, homomorphisms}
A \textbf{monoid} is a set with an associative product and an identity. Its invertible elements form its \textbf{group
of units}. The submonoid \textbf{generated} by a subset $S$ consists of all finite products $s_1\cdots s_k$ of elements
of $S$, the empty product being the identity. A monoid acts on a set when $1\cdot x=x$ and $s\cdot(t\cdot x)=(st)\cdot
x$, and the \textbf{monoid orbit} of $x$ is $\{s\cdot x\}$. A \textbf{homomorphism} of monoids preserves products and
the identity. Examples are $\mathrm{Diag}_m$, $M_2(\R)$ and $\operatorname{End}(\Theta)$ under multiplication, with
units $T_m$, $\GL_2(\R)$ and $\GL(\Theta)$, and finite multisets under disjoint union $\uplus$. Additive monoids
$\mathrm{Mon}(C)$ are in \bgref{B.14}.
\end{bgentry}

Merge-and-restart reaches the orbit of the monoid that one cycle generates, which is $\theta_0+\mathrm{Mon}(C)$ for an
additive method $\rho=\theta_0+\delta$ with shape $C=\Im\delta$ (\bgref{B.14}) and $\langle S\rangle\cdot\theta_0$ for a
method of action type (\bgref{D.12}) whose cycle applies elements of $S$. RoAd$_2$'s blocks range over $M_2(\R)$. For a
fixed query, attention's pair (total weight, weighted sum of values) is a homomorphism from multisets of key--value
pairs to $(\R_{\ge0}\times\R^{d_v},+)$ (\bgref{D.34}).

\fbAusedin{\thmref{Theorem}{V.3}, restarting (\S\ref{sec:base-rebase}), DoRA (\S\ref{sec:erlangen-dora}),
  \thmref{Proposition}{II.9}, \thmref{Theorem}{VI.5}, attention (\S\ref{sec:merge-attention}).}

\begin{bgfact}{D.11}{Cartan--Dieudonné theorem; Scherk's sharp form}
For nonzero $u\in\R^n$, the \textbf{reflection} $H_u=I-2uu\T/\|u\|^2$ fixes the hyperplane $u^\perp$ and has $\det
H_u=-1$. Every $R\in\Ort(n)$ is a product of at most $n$ reflections (Cartan--Dieudonné). In Scherk's sharp form, if
$k=\rk(R-I)$, then $R$ is a product of exactly $k$ reflections and of no fewer. Hence $\det R=(-1)^k$, and every
$R\in\SO(n)$ is a product of $\rk(R-I)$ reflections, an even number. See \citet{xscherk1950decomposition} and
\citet[Ch.~III]{xartin1957geometric}.
\end{bgfact}

\paragraph{The idea of the proof}
A product of $j$ reflections fixes an intersection of $j$ hyperplanes, so $\rk(R-I)\le j$. Conversely, pick $x$ with
$Rx\ne x$ and put $u=Rx-x$. Then $H_u$ swaps $Rx$ and $x$ and fixes every fixed vector of $R$, so $H_uR$ fixes one more
dimension, and induction on $k$ finishes.

HRA multiplies $r$ reflections, so $\rk(R-I)\le r$, and $\det R=1$ for even $r$. Scherk's form gives the converse. For
injective $W_0$ and even $r$, HRA reaches all of $(W_0+\Mr)\cap W_0\cdot\SO(n)$, padding with pairs $H_uH_u=I$, and
re-merged HRA reaches $W_0\cdot\SO(n)$ within $\lceil n/r\rceil$ cycles.

\fbAusedin{\thmref{Theorem}{II.7}, HRA is a meet (\S\ref{sec:erlangen-hra}), the corollaries of \thmref{Theorem}{V.3}
  (\S\ref{sec:base-rebase}).}

% ----------------------------------------------------------------------------------------------------------------------
\subsubsection*{Actions, orbits and invariants}

\begin{bgentry}{D.12}{group actions}
A (left) \textbf{action} of a group $G$ on a set $X$ is a map $(g,x)\mapsto g\cdot x$ with $e\cdot x=x$ and
$g\cdot(h\cdot x)=(gh)\cdot x$. A \textbf{right} action $x\cdot g$ satisfies $x\cdot(gh)=(x\cdot g)\cdot h$, and $g\cdot
x:=x\cdot g^{-1}$ turns it into a left one. The action is \textbf{smooth} if $G\times X\to X$ is smooth, and
\textbf{linear} if $X$ is a vector space and each map $x\mapsto g\cdot x$ is linear.
\end{bgentry}

On weights, $W\mapsto WR$ is a right action of $\Ort(n)$ on the input side and $W\mapsto PW$ a left action of $\Ort(m)$
on the output side, and $(g,h)\cdot W=gWh^{-1}$ is a left action of $\GL_m\times\GL_n$. On LoRA's factors,
$g\cdot(B,A)=(Bg^{-1},gA)$ is a linear left action of $\GL_r$. A method is of \textbf{action type} when
$\rho(q)=\gamma(q)\cdot\theta_0$ for a map $\gamma$ into a group acting on $\Theta$, as for OFT and HRA. (This $\gamma$
follows the exposés and is unrelated to the cotangent signal $\gamma(t)$ of \bgref{C.2} and \bgref{D.26}.) (IA)$^3$ has
the same form with the monoid $\mathrm{Diag}_m$ in place of a group, and DoRA's magnitude applies that monoid to LoRA's
weight.

\fbAusedin{\thmref{Definition}{I.2}, the notation (Table~\ref{tab:notation}), Erlangen (\S\ref{sec:erlangen}),
  \thmref{Theorem}{II.5}, \thmref{Theorem}{III.5}.}

\begin{bgentry}{D.13}{orbits, stabilisers, free actions}
The \textbf{orbit} of $x$ is $G\cdot x=\{g\cdot x:g\in G\}$, its closure is the \textbf{orbit closure}, and its
\textbf{stabiliser} is the closed subgroup $G_x=\{g:g\cdot x=x\}$. For a smooth action of a Lie group, $g\mapsto g\cdot
x$ induces an injective immersion $G/G_x\to X$ with image $G\cdot x$, so the orbit is an immersed submanifold of
dimension $\dim G-\dim G_x$. The \textbf{infinitesimal stabiliser} $\{\xi\in\mathfrak g:\xi\cdot x=0\}$ is the Lie
algebra of $G_x$. The action is \textbf{free} if every stabiliser is trivial, and the \textbf{generic stabiliser} is the
stabiliser at a generic point, one in a suitable dense open set. Each map $x\mapsto g\cdot x$ carries an orbit onto
itself, so all its points look alike; orbits are \textbf{homogeneous}.
\end{bgentry}

$\GL_r$ acts freely on pairs with $\rk A=r$, since $gA=A$ forces $g=I$, so LoRA's generic fibre, a single orbit, has
dimension $r^2$. LoHa's joint gauge $\GL_{r_1}\times\GL_{r_2}\times T_m\times T_n$ has a one-dimensional generic
stabiliser, so its orbits have dimension $r_1^2+r_2^2+m+n-1$. By homogeneity, the two-sided torus orbit
$\{\diag(a)W_0\diag(b)\}$ has dimension $m+n-c(W_0)$ at each of its points, with $c(W_0)$ the number of connected
components of the support graph, the bipartite graph on the $m$ rows and $n$ columns with an edge $\{i,j\}$ where
$(W_0)_{ij}\ne0$. The images of (IA)$^3$ and RoAd$_1$ are orbit closures.

\fbAusedin{Erlangen (\S\ref{sec:erlangen}), \thmref{Theorem}{II.5}, \thmref{Proposition}{II.9}, the image germ
  (Exposé~\ref{sec:image}), \thmref{Proposition}{II.12}, Rosetta entries 6 and 11 (Exposé~\ref{sec:rosetta}).
  \textsf{Reading:} \citet[Ch.~7 and~21]{xlee2013introduction}.}

\begin{bgentry}{D.14}{Haar measure; invariant measures on homogeneous spaces}
A compact Lie group $K$ carries a unique invariant probability measure, its \textbf{Haar measure}. If $K$ acts
transitively on $X$, so that $X\cong K/K_x$ is a \textbf{homogeneous space}, then $X$ carries a unique $K$-invariant
probability measure, the image of Haar measure under $k\mapsto k\cdot x$. The Grassmannian $\Gr(r,n)$ of $r$-dimensional
subspaces of $\R^n$ (\bgref{C.15}) is a homogeneous space of $\Ort(n)$.
\end{bgentry}

If $A_0\in\R^{r\times n}$ has i.i.d.\ Gaussian entries, its law is unchanged by $A_0\mapsto A_0R$ for $R\in\Ort(n)$, and
$A_0$ has rank $r$ almost surely. So $\operatorname{row}A_0$ is an $\Ort(n)$-invariant random point of $\Gr(r,n)$, and
by uniqueness it is Haar-distributed. Vanilla LoRA with a Gaussian $A_0$ therefore picks a Haar-random point of
$\Gr(r,n)$, the space that classifies its stratum $\mathrm T_A$. Kaiming-uniform entries are not rotation invariant, so
under HF's default the point is not exactly Haar-distributed.

\fbAusedin{the three strata (\S\ref{sec:strata-three}), \thmref{Theorem}{II.2}. \textsf{Reading:}
  \citet[Ch.~3]{xmattila1995geometry}.}

\begin{bgentry}{D.15}{invariants and complete invariants}
A map $f:X\to Y$ is \textbf{invariant} under an action if $f(g\cdot x)=f(x)$ for all $g$ and $x$, so it is constant on
orbits. It is a \textbf{complete invariant} if moreover $f(x)=f(x')$ implies $x'\in G\cdot x$; it then identifies the
orbit space $X/G$ with the image of $f$.

\smallskip
{\centering\small\renewcommand{\arraystretch}{1.12}
\begin{tabularx}{\linewidth}{@{}XX@{}}
  \toprule
  \headrow \textbf{action on $W\in\R^{m\times n}$} & \textbf{complete invariant}\\
  \midrule
  $W\mapsto WR$, $R\in\Ort(n)$ (input side) & $K_{\mathrm{out}}=WW\T$\\
  $W\mapsto PW$, $P\in\Ort(m)$ (output side) & $K_{\mathrm{in}}=W\T W$\\
  $W\mapsto PWR$, $P\in\Ort(m)$, $R\in\Ort(n)$ & singular values with multiplicity\\
  $W\mapsto DW$, $D\in T_m^+$, $W$ without zero rows & the rays $\R_{>0}w_i$ of the rows\\
  $W\mapsto E\odot W$, $E\in(\R^\times)^{m\times n}$ & the zero pattern\\
  \bottomrule
\end{tabularx}\par}
\end{bgentry}

If a method's image lies in one orbit, every invariant of that action stays fixed through training and merging. This is
why exactly orthogonal methods never change $K_{\mathrm{out}}$ or the spectrum (\thmref{Corollary}{II.6}). Isomorphism
invariants of methods, such as $S_1$, $d(M)$ and the first-order deficit, play the same role one level up, since methods
with different values cannot be isomorphic.

\fbAusedin{\thmref{Theorem}{II.5}, complete invariants and the Gram lab (\S\ref{sec:erlangen-invariants}),
  \thmref{Corollary}{II.6}, \thmref{Observation}{I.6}, the seven borrowed words (Exposé~\ref{sec:analogy}).}

\begin{bgfact}{D.16}{first fundamental theorem for $\Ort(n)$; Witt's extension theorem}
(First fundamental theorem.) Every polynomial in the rows $w_1,\dots,w_m\in\R^n$ of $W$ that is invariant under
$W\mapsto WR$, $R\in\Ort(n)$, is a polynomial in the inner products $\langle w_i,w_j\rangle$, the entries of $WW\T$. For
a compact group acting linearly, invariant polynomials separate orbits, so $WW\T$ is a complete invariant. (Witt's
extension theorem, Euclidean case.) A linear isometry between two subspaces of $\R^n$ extends to an element of
$\Ort(n)$. See \citet[Ch.~II]{xweyl1939classical} and \citet[Ch.~5]{xgoodman2009symmetry}.
\end{bgfact}

\paragraph{Witt gives the orbit separation directly}
If $WW\T=W'{W'}\T$, the map $\sum_ic_iw_i\mapsto\sum_ic_iw_i'$ is well defined and isometric on the row spaces, because
$\|\sum_ic_iw_i\|^2=c\T WW\T c$. Extend it by an isometry between the orthogonal complements to $R\T\in\Ort(n)$; then
$W'=WR$.

This is item (1) of \thmref{Theorem}{II.5}, the reason exactly orthogonal input-side methods are pinned to
$K_{\mathrm{out}}(W_0)$. Applied to $W\T$, it gives $K_{\mathrm{in}}$ for the output side.

\fbAusedin{\thmref{Theorem}{II.5}, the known results (Exposé~\ref{sec:analogy}).}

\begin{bgentry}{D.17}{Klein's Erlangen programme}
\citet{xklein1872vergleichende} proposed to define a geometry by a space together with a group acting on it, and to
study the properties the group preserves. Euclidean geometry keeps lengths and angles, the invariants of rigid motions.
Affine geometry keeps parallelism and ratios along parallel lines but not lengths or angles, because its group is
larger. A larger group leaves fewer invariants.
\end{bgentry}

For a group-type method the reading is a theorem, because the image lies in one orbit and so the method can change
nothing the group preserves. For a cone such as LoRA's image the word only names its \textbf{covariance group}, the
largest $\mathcal G_M\subseteq\GL_m\times\GL_n$ with $\Im M(gW_0h^{-1})=g\,\Im M(W_0)\,h^{-1}$ (\bgref{D.18}). For LoRA
this is all of $\GL_m\times\GL_n$, for (IA)$^3$ on keys and values it is $\mathrm{Mon}_m\times\GL_n$, and for the PiSSA
split it is $\CO(m)\times\CO(n)$.

\fbAusedin{Erlangen and what ``Erlangen'' means here (\S\ref{sec:erlangen}), \thmref{Definition}{II.14}, the borrowed
  words (Exposé~\ref{sec:rosetta}), the seven borrowed words (Exposé~\ref{sec:analogy}). \textsf{Reading:}
  \citet[Ch.~4]{xsharpe1997differential}.}

\begin{bgentry}{D.18}{equivariance and invariance of maps, update rules and merge rules}
A map $F:X\to Y$ between spaces with $G$-actions is \textbf{equivariant} if $F(g\cdot x)=g\cdot F(x)$, and
\textbf{invariant} if $F(g\cdot x)=F(x)$. An update rule on parameters is equivariant under a subgroup $H$ of the gauge
if, started from $h\cdot q$ with correspondingly transformed optimizer state, it produces $h$ applied to the iterates
started from $q$. A rule that uses randomness is \textbf{equivariant in distribution} if $F(g\cdot x)$ and $g\cdot F(x)$
have the same law for every $x$.
\end{bgentry}

On LoRA's full-rank locus, with channel-uniform hyperparameters, the largest such subgroup of $\GL_r$ is $\Ort(r)$ for
SGD, $B_r$ for Adam and all of $\GL_r$ for undamped scaled GD (\bgref{D.23}, \thmref{Theorem}{III.12}). An equivariant
rule started from gauge-related points traces the same weights, since $\rho(h\cdot q)=\rho(q)$. Merge rules should be
invariant under each adapter's gauge. Task arithmetic \citep{ilharco2022editing} $\sum_iw_is_iB_iA_i$ is, while
LoraHub's \citep{huang2023lorahub} $(\sum_iw_iB_i)(\sum_jw_jA_j)$ is not. On a weight space $\R^N$, DARE's
\citep{yu2023language} random drop-and-rescale is equivariant under the invertible diagonal matrices $T_N$ for each
fixed mask, and under the monomial matrices $\mathrm{Mon}_N$ (\bgref{D.3}) in distribution.

\fbAusedin{\thmref{Definition}{II.14}, \thmref{Observation}{II.15}, \thmref{Theorem}{III.12},
  \thmref{Corollary}{III.13}, Figure~\ref{fig:gauge}, \thmref{Proposition}{V.4}.}

\begin{bgentry}{D.19}{multihomogeneous maps; multidegree}
Split the parameters into blocks $q=(q_1,\dots,q_k)$. A map $\delta$ is \textbf{multihomogeneous} of
\textbf{multidegree} $(d_1,\dots,d_k)$ if $\delta(c_1q_1,\dots,c_kq_k)=c_1^{d_1}\cdots c_k^{d_k}\,\delta(q)$ for all
real $c_i$. LoRA's $BA$ has multidegree $(1,1)$, and LoHa's $(B_1A_1)\odot(B_2A_2)$ has multidegree $(1,1,1,1)$.
\end{bgentry}

If some $d_i\ge1$, the image of $\delta$ contains $0$ and is closed under multiplication by every nonnegative scalar, so
it is a cone with vertex $0$ (\bgref{C.10}). For a method $\rho=\theta_b+s\,\delta$ started where $\delta(q_0)=0$, so
that $\theta_b=\theta_0$, the image is the cone $\theta_0+s\,\Im\delta$ with vertex $\theta_0$. For $c\in\R^k_{>0}$, the
change $q_i\mapsto c_iq_i$ together with $s\mapsto s\prod_ic_i^{-d_i}$ leaves $\rho$ unchanged. That rescaling is an
isomorphism of methods and the source of the hyperparameter torus (\bgref{D.20}).

\fbAusedin{the image germ (Exposé~\ref{sec:image}), \thmref{Theorem}{III.10}, the seven axes (\S\ref{sec:atlas-axes}).}

\begin{bgentry}{D.20}{orbit space of a positive torus}
Let $\R^k_{>0}$ act on positive variables $x\in\R^N_{>0}$ by $x_j\mapsto x_j\prod_ic_i^{a_{ji}}$, with a real weight
matrix $a=(a_{ji})$. In logarithms this is translation by the column space of $a$, so the orbits are parallel affine
subspaces and the orbit space is a vector space of dimension $N-\rk a$. The monomials $\prod_jx_j^{b_j}$ with $b\T a=0$
are invariant, and a basis of such vectors $b$ gives a complete set of invariants.
\end{bgentry}

Under the hypotheses of \thmref{Theorem}{III.10}, taken with Adam at $\varepsilon=0$, no weight decay or clipping and a
shared schedule (otherwise the per-block $\varepsilon$, decay and clipping thresholds rescale as well), the torus acts
on LoRA's hyperparameters by $s\mapsto s/(c_Ac_B)$, $\eta_A\mapsto c_A\eta_A$, $\eta_B\mapsto c_B\eta_B$,
$\sigma_A\mapsto c_A\sigma_A$, with $\sigma_A$ the scale of $A$'s initialisation, without changing the weight
trajectory. The invariants are $s\eta_A\eta_B$ and $\sigma_A/\eta_A$, so only two of the four knobs matter. LoRA+
\citep{hayou2024lora} with ratio $\lambda$ replaces $\eta_B$ by $\lambda\eta_B$, which changes the invariants exactly as
replacing $\alpha$ by $\lambda\alpha$ does. Under SGD the rates scale by $c_i^2$, and the invariants become
$s^2\eta_A\eta_B$ and $\sigma_A^2/\eta_A$.

\fbAusedin{\thmref{Theorem}{III.10}, \thmref{Corollary}{III.11}, Figure~\ref{fig:orbit}, Rosetta entry 23
  (Exposé~\ref{sec:rosetta}), \thmref{Prediction}{X.1}.}

% ----------------------------------------------------------------------------------------------------------------------
\subsubsection*{Gauge, quotients and metrics}

\begin{bgentry}{D.21}{gauge symmetry of a parametrisation}
The \textbf{gauge group} of a method $\rho:Q\to\Theta$ is the group of diffeomorphisms $\varphi$ of $Q$ with
$\rho\circ\varphi=\rho$. If a Lie group $G$ acts on $Q$ with $\rho(g\cdot q)=\rho(q)$, each orbit lies inside a fibre
$\rho^{-1}(\theta)$ (\bgref{C.6}), so $d(M)\le\lvert M\rvert-\dim(\text{generic orbit})$. The action is \textbf{locally
transitive} on a fibre when its orbits are open in the fibre, that is, when orbits and fibre have the same dimension
there; on generic fibres this gives $d(M)=\lvert M\rvert-\dim(\text{generic orbit})$, and we then call $G$ the generic
gauge. Here the word means only this symmetry group; no connection is involved.
\end{bgentry}

LoRA's $\GL_r$ acts freely and locally transitively on full-rank fibres, so $d=r(m+n)-r^2=r(m+n-r)$, and $r^2$
parameters are gauge waste. HRA's generic fibres have dimension $r(r+1)/2$, while its $r$ scalings $u_i\mapsto c_iu_i$
sweep out only $r$ of those dimensions, so the scalings do not fill the fibres; \thmref{Theorem}{II.7}(a) describes the
rest of the gauge as non-linear. For $r=2$ the fibre through $(u_1,u_2)$ has dimension $3$, made up of the two scalings
and of rotating $u_1$ and $u_2$ by a common angle inside their plane, which keeps the angle between them and hence
$H_{u_1}H_{u_2}$, the rotation of that plane by twice the angle. The plane depends on the point, so this move is not a
linear action on parameters. The \textbf{Hurwitz moves} $(u_i,u_{i+1})\mapsto(H_{u_i}u_{i+1},u_i)$ preserve the product
of reflections, since $H_aH_b=H_{H_ab}H_a$, and they satisfy the braid relations, so the braid group on $r$ strands,
unrelated to the signed permutations $B_r$, acts on HRA's parameters. These moves are discrete and add no dimension to
the fibre.

\fbAusedin{\thmref{Definition}{I.5}, invariants of a method (\S\ref{sec:slice-invariants}), the notation
  (Table~\ref{tab:notation}), rank and dimension (\S\ref{sec:cut}), \thmref{Proposition}{II.12}, \thmref{Theorem}{II.7},
  the analogy ``Noether; gauge'' (\S\ref{sec:fibres-charges}), the seven borrowed words (Exposé~\ref{sec:analogy}).}

\begin{bgfact}{D.22}{quotients and descent; the quotient manifold theorem}
For an action of $G$ on $X$, the \textbf{orbit space} $X/G$ is the set of orbits, with the map $\pi:x\mapsto G\cdot x$.
A map $f:X\to Y$ that is constant on orbits factors uniquely as $f=\bar f\circ\pi$; this is \textbf{descent}. (Quotient
manifold theorem.) If a Lie group acts smoothly, freely and properly on a manifold $X$, then $X/G$ is a manifold of
dimension $\dim X-\dim G$, $\pi$ is a submersion, and $\bar f$ is smooth whenever $f$ is. Properness means that
$(g,x)\mapsto(g\cdot x,x)$ is a proper map, one under which preimages of compact sets are compact
\citep[Thm~21.10]{xlee2013introduction}.
\end{bgfact}

$\GL_r$ acts freely and properly on pairs $(B,A)$ of full rank, and $(B,A)\mapsto BA$ identifies the quotient with
$\MM_r$, of dimension $r(m+n)-r^2$. A rule for combining LoRA factors is a function of the adapters exactly when it is
constant on the fibres of $(B_i,A_i)\mapsto(B_iA_i)$, and a cometric on factors that depends only on $W$ descends to
$\MM_r$.

\fbAusedin{\thmref{Proposition}{V.4}, \thmref{Proposition}{III.4}, \thmref{Observation}{II.15}, when a change of chart
  does not matter (\S\ref{sec:fibres-naturality}).}

\begin{bgentry}{D.23}{metrics, cometrics and the pulled-back cometric}
A \textbf{metric} on a vector space $V$ is a positive-definite inner product, equivalently a symmetric positive-definite
map $g:V\to V^*$, that is, $\langle gv,v\rangle>0$ for $v\ne0$. A \textbf{cometric} is a symmetric positive semidefinite
map $K:V^*\to V$. It turns a covector, such as a gradient $\nabla_\theta L$, into a velocity $-K\nabla_\theta L$. A
metric gives the cometric $K=g^{-1}$, while a cometric may be degenerate. Metrics pull back along immersions, and
cometrics push forward along any smooth map. For instance, gradient flow $\dot q=-\Lambda\nabla_q(L\circ\rho)$ moves the
weights by $\dot\theta=-K_q\nabla_\theta L$, with
\[K_q=D\rho_q\,\Lambda\,D\rho_q\T.\]
\end{bgentry}

For LoRA with \textbf{block-scalar rates}, one rate per factor, $\Lambda=\diag(\eta_BI,\eta_AI)$ on $(B,A)$ and
$K(G)=s^2(\eta_BGA\T A+\eta_ABB\T G)$. Two points of one fibre give the same weight velocity for every loss exactly when
their cometrics agree. At points where $A$ has full row rank and $r<n$, this holds for LoRA's linear gauge and plain
gradient descent exactly along $\Ort(r)$. A rule descends to a cometric on weights when it gives the same weight
velocity at all points of each fibre; with a point-dependent preconditioner in place of $\Lambda$, undamped scaled GD
\citep{xtong2021accelerating,zhang2024riemannian} and the SGD variant of LoRA-Pro \citep{wang2024lorac} do, on the
full-rank locus. Scaled GD multiplies the $B$-gradient by $(AA\T)^{-1}$ and the $A$-gradient by $(B\T B)^{-1}$, and
undamped means that no ridge term $\delta I$ is added before inverting (\thmref{Proposition}{III.4}).

\fbAusedin{\thmref{Definition}{III.1}, \thmref{Observation}{III.2}, \thmref{Theorem}{III.3},
  \thmref{Proposition}{III.4}, the metric a method carries (\S\ref{sec:fibres-cometric}), the three invariants
  (Exposé~\ref{sec:thesis}).}

\begin{bgentry}{D.24}{Riemannian gradient flow; tangent projection; retraction}
On a submanifold $N\subseteq\Theta$ with the metric induced by the Frobenius inner product, the Riemannian gradient of
$L|_N$ at $\theta$ is the orthogonal projection $\Pi_{T_\theta N}\nabla L(\theta)$ of the ambient gradient onto the
tangent space, and the \textbf{Riemannian gradient flow} is $\dot\theta=-\Pi_{T_\theta N}\nabla L(\theta)$. A discrete
step along the tangent space leaves $N$, so it is followed by a \textbf{retraction}, a smooth map $R_\theta$ from a
neighbourhood of $0$ in $T_\theta N$ to $N$ with $R_\theta(0)=\theta$ and $DR_\theta(0)=\id$.
\end{bgentry}

At $W=W_0+X$ with $X\in\MM_r$, $U=\operatorname{col}X$ and $V=\operatorname{row}X$, the tangent space is
$\{D:P_{U^\perp}DP_{V^\perp}=0\}$ (\bgref{C.12}) and the projection is $P_UG+GP_V-P_UGP_V$. This is LoRA-Pro's
equivalent gradient, so the SGD variant of LoRA-Pro follows the Riemannian gradient flow on $W_0+\MM_r$. Truncating the
SVD of $X+\xi$ to rank $r$ is a retraction onto $\MM_r$, the one Riemannion \citep{bogachev2025lora} uses.

\fbAusedin{\thmref{Proposition}{III.4}, when a change of chart does not matter (\S\ref{sec:fibres-naturality}),
  \thmref{Theorem}{III.12}. \textsf{Reading:} \citet[Ch.~3 and~4]{xabsil2008optimization}.}

\begin{bgfact}{D.25}{distributions, Lie brackets and Frobenius' theorem}
A \textbf{distribution} $D$ of rank $k$ on an open set $U$ assigns to each point $\theta$ a $k$-dimensional subspace
$D_\theta$ of directions, smoothly in $\theta$; the word is unrelated to probability distributions. The \textbf{Lie
bracket} of vector fields is $[X,Y]=DY\cdot X-DX\cdot Y$; flowing along $X$, $Y$, $-X$, $-Y$ for time $t$ misses the
starting point by $t^2[X,Y]+O(t^3)$. $D$ is \textbf{involutive} if $[X,Y]$ lies in $D$ whenever $X$ and $Y$ do.
(Frobenius.) $D$ is involutive iff every point lies on a $k$-dimensional \textbf{integral manifold}, a submanifold whose
tangent space at each of its points $\theta$ is $D_\theta$. A maximal connected integral manifold is a \textbf{leaf}
\citep[Ch.~19]{xlee2013introduction}. The Frobenius algebras of \bgref{B.6} and the Frobenius norm share only the name.
\end{bgfact}

On $\R^3$, $D=\spanop\{\partial_x,\ \partial_y+x\,\partial_z\}$ is not involutive, since
$[\partial_x,\partial_y+x\,\partial_z]=\partial_z\notin D$, so no surface is tangent to $D$. A backward method with
anchor $a(\theta)$ chooses $D_\theta=\operatorname{im}a(\theta)$, and it reads as a forward method only along integral
manifolds. GaLore \citep{zhao2024galore} with a projector fixed on each period, LISA \citep{pan2024lisa} and MeZO
\citep{malladi2023fine} have constant distributions, which are involutive with affine leaves. Full-batch GaLore with
$T=1$, read as $D_W=\operatorname{Hom}(\R^n,U_r(\nabla L(W)))$ wherever a singular-value gap exists, already fails to be
involutive for $L=\tfrac12\|W\|_F^2$.

\fbAusedin{when a backward method is a forward one (\S\ref{sec:lens-integrability}), \thmref{Proposition}{IV.3}.}

% ----------------------------------------------------------------------------------------------------------------------
\subsubsection*{Conservation and balance}

\begin{bgfact}{D.26}{Noether's theorem for gradient flows}
Let a Lie group act linearly on $Q=\R^N$ with $\rho(k\cdot q)=\rho(q)$, let $\Lambda\succ0$ be a constant rate matrix,
and let $\gamma(t)$ be any time-dependent covector on $\Theta$, a gradient or not. Along every solution of $\dot
q=-\Lambda\,D\rho_q\T\gamma(t)$, and for every generator $\xi$ of the action with $\Lambda^{-1}\xi$ symmetric, the
\textbf{charge}
\[J_\xi(q)=\tfrac12\langle q,\Lambda^{-1}\xi\,q\rangle\]
is constant \citep{kunin2020neural,zhao2022symmetries,marcotte2023abide}.
\end{bgfact}

\paragraph{The one-line proof}
Invariance of $\rho$ gives $D\rho_q(\xi q)=0$. Hence $\tfrac{d}{dt}J_\xi=\langle\Lambda^{-1}\xi q,\dot
q\rangle=-\langle\xi q,D\rho_q\T\gamma\rangle=-\langle D\rho_q\,\xi q,\gamma\rangle=0$.

For scalar LoRA $w=w_0+ba$ with rates $\eta_B,\eta_A$, the charge $b^2/\eta_B-a^2/\eta_A$ is conserved, so every
trajectory runs along a hyperbola, the Noether hyperbolas of Figure~\ref{fig:noether}(a). \citet{xnoether1918invariante}
derived conservation laws from symmetries of a Lagrangian; this version needs no Lagrangian and holds for any loss and
data. For LoRA with rates $\eta_B,\eta_A$, the generators $\xi_X(B,A)=(BX,-XA)$ pass the test exactly for symmetric $X$,
and the $r(r+1)/2$ charges are the entries of $\Phi=B\T B/\eta_B-AA\T/\eta_A$. Minibatches, LoRA dropout, global-norm
clipping and DoRA's detached norm keep the required form. Momentum and Adam do not, and under discrete SGD the charge
drifts, by an amount relative to the charge of order $\eta^2$ per step.

\fbAusedin{\thmref{Theorem}{III.5}, isometries and charges and the analogy ``Noether; gauge''
  (\S\ref{sec:fibres-charges}), Figure~\ref{fig:noether}, the things one can check (Appendix~\ref{app:repro}),
  \thmref{Prediction}{X.8}, the seven borrowed words (Exposé~\ref{sec:analogy}).}

\begin{bgentry}{D.27}{Cartan decomposition; polar decomposition}
The map $\vartheta(\xi)=-\xi\T$ is an automorphism of $\gl_r$ of order two, a \textbf{Cartan involution}. Its $+1$ and
$-1$ eigenspaces give the \textbf{Cartan decomposition} $\gl_r=\so(r)\oplus\Sym_r$,
$\xi=\tfrac12(\xi-\xi\T)+\tfrac12(\xi+\xi\T)$, with $[\so,\so]\subseteq\so$, $[\so,\Sym]\subseteq\Sym$ and
$[\Sym,\Sym]\subseteq\so$. Its group version is the \textbf{polar decomposition}, which says that $(k,S)\mapsto
k\,e^{S}$ is a diffeomorphism $\Ort(r)\times\Sym_r\to\GL_r$. Hence $g\mapsto g\T g$ identifies the cosets $\Ort(r)g$
with the positive-definite matrices, $\Ort(r)\backslash\GL_r\cong\Sym^+_r$.
\end{bgentry}

With block-scalar rates (\bgref{D.23}), the antisymmetric generators $X\in\so(r)$ of LoRA's gauge act by rotations
$(BO\T,OA)$ that preserve the rate metric, so the optimizer cannot feel them; they are \textbf{isometries}. The
symmetric generators $X\in\Sym_r$ give the conserved charges of \bgref{D.26}. For general rates the split needs
$\mathfrak g$ to be closed under $\xi\mapsto\Lambda\xi\T\Lambda^{-1}$. Gradient descent on LoRA is thus a family of
dynamics indexed by $\Ort(r)\backslash\GL_r$, and the Kempf--Ness argument needs only curves $e^{t\xi}$ with $\xi$
symmetric.

\fbAusedin{\thmref{Theorem}{III.5}, isometries and charges (\S\ref{sec:fibres-charges}), the three invariants and the
  thesis (Exposé~\ref{sec:thesis}), \thmref{Proposition}{III.7}, when a change of chart does not matter
  (\S\ref{sec:fibres-naturality}). \textsf{Reading:} \citet[\S VI.2]{xknapp2002lie}; \citet[\S 7.3]{xhorn2013matrix}.}

\begin{bgentry}{D.28}{balancedness; the conserved $B\T B-AA\T$}
A factorisation $W=BA$ is \textbf{balanced} if $B\T B=AA\T$, and rate-weighted balanced if $\Phi=B\T
B/\eta_B-AA\T/\eta_A=0$, that is, $B\T B=cAA\T$ with $c=\eta_B/\eta_A$. For a deep linear network $W=W_L\cdots W_1$ the
same word means $W_{j+1}\T W_{j+1}=W_jW_j\T$ for every $j$.
\end{bgentry}

Under gradient flow with block-scalar rates, $\Phi$ is conserved (\bgref{D.26}), so rate-weighted balancedness, once it
holds, holds forever; at equal rates the two notions agree. Zero-product starts $B_0A_0=0$ have $\langle B_0\T
B_0,A_0A_0\T\rangle_F=0$, which makes them maximally unbalanced for given norms of the two Gram matrices, since the
inner product of two positive semidefinite matrices is never negative. From $B_0=0$, conservation gives $B\T
B=\tfrac{\eta_B}{\eta_A}(AA\T-A_0A_0\T)$, so $AA\T$ never drops below $A_0A_0\T$ in the Loewner order. HF's PiSSA starts
with $B_0\T B_0=A_0A_0\T$. When $\Phi$ is isotropic, singular values of $\Delta W$ well above the crossover scale
$\sigma^*$ of \thmref{Theorem}{III.9}, set by the charge, the scale and the rates, follow the balanced dynamics of
\citet{arora2018optimization}.

\fbAusedin{isometries and charges (\S\ref{sec:fibres-charges}), \thmref{Proposition}{III.7},
  \thmref{Proposition}{III.8}, \thmref{Theorem}{III.9}, the bibliography. \textsf{Reading:}
  \citet{arora2018optimization}; \citet{du2018algorithmic}.}

\begin{bgentry}{D.29}{the moment map}
Let a group $G\subseteq\GL(V)$ that is closed under transpose act on a real inner-product space $V$, and split
$\mathfrak g=\mathfrak k\oplus\mathfrak p$ into its antisymmetric and symmetric elements. The (real) \textbf{moment map}
$\mu:V\to\mathfrak p$ is defined by
\[\langle\mu(v),\xi\rangle=\tfrac12\tfrac{d}{dt}\Big|_{t=0}\|e^{t\xi}v\|^2=\langle\xi v,v\rangle\qquad(\xi\in\mathfrak p).\]
It records how the norm changes along the symmetric directions of the orbit; the antisymmetric directions preserve the
norm.
\end{bgentry}

For LoRA's gauge, $\tfrac{d}{dt}\big|_{0}\big(\|Be^{-tX}\|^2+\|e^{tX}A\|^2\big)=-2\tr(X\mu)$ for symmetric $X$, with
$\mu=B\T B-AA\T$, so this $\mu$ is the moment map up to sign. Its zero set is the balanced slice, and at equal rates
$\eta_A=\eta_B=\eta$ it equals $\eta\Phi$, a multiple of the Noether charge.

\fbAusedin{the balanced slice (\S\ref{sec:fibres-balance}), \thmref{Proposition}{III.7}, \thmref{Corollary}{III.6},
  Figure~\ref{fig:noether}.}

\begin{bgfact}{D.30}{Kempf--Ness theorem, real form}
Let $G\subseteq\GL(V)$ be a real algebraic group, cut out by polynomial equations, that is closed under transpose (such
a group is \textbf{reductive}), for instance $\GL_r$ acting on LoRA factors, and let $K=G\cap\Ort(V)$. For $v\in V$, (i)
$g=e$ is a critical point of $g\mapsto\|g\cdot v\|^2$ iff $\mu(v)=0$, and then $v$ has least norm on $G\cdot v$; (ii)
$G\cdot v$ is closed iff it contains a point with $\mu=0$; (iii) on a closed orbit, the points with $\mu=0$ form a
single $K$-orbit; (iv) the closure of every orbit contains exactly one closed orbit
\citep{xkempf1979length,xbirkes1971orbits,xrichardson1990minimum}.
\end{bgfact}

\paragraph{The idea of the proof}
For symmetric $\xi$, $\|e^{t\xi}v\|^2=\sum_ie^{2t\lambda_i}\|v_i\|^2$ in an eigenbasis of $\xi$, which is convex in $t$.
By the polar decomposition every orbit point is $k\,e^{\xi}v$ with $k\in K$ norm-preserving, so critical points are
minima and two minima differ by an element of $K$.

For $m=n=r=1$, $t\in\R^\times$ acts by $t\cdot(b,a)=(b/t,ta)$ and $\mu=b^2-a^2$. The orbits are the hyperbolas
$ba=c\ne0$, which are closed and meet $\mu=0$ in the two points $\pm(b,a)$ with $\lvert b\rvert=\lvert a\rvert$, one
$\Ort(1)$-orbit; the two punctured axes, which are not closed and avoid $\mu=0$; and the origin. For LoRA in general,
$K=\Ort(r)$, and a full-rank orbit is closed, being a whole fibre of $(B,A)\mapsto BA$. The least-norm factorisations of
a fixed $BA$ are the balanced ones, they form one $\Ort(r)$-orbit, and there $\|B\|_F^2+\|A\|_F^2=2\|BA\|_*$. Weight
decay adds a norm to the loss (\thmref{Corollary}{III.6}). With learning-rate-coupled decay, which multiplies the decay
of each factor by its rate, $\dot\Phi=-2\lambda\mu$, so the equilibria lie on $\mu=0$, and $\mu$ decays exponentially at
equal rates. With per-time decay, the same $\lambda$ for both factors, $\dot\Phi=-2\lambda\Phi$, so $\Phi\to0$.

\fbAusedin{the balanced slice (\S\ref{sec:fibres-balance}), \thmref{Proposition}{III.7}, \thmref{Corollary}{III.6},
  combining (\S\ref{sec:base-combine}), the known results (Exposé~\ref{sec:analogy}), the classification
  (Exposé~\ref{sec:atlas}).}

\begin{bgfact}{D.31}{the closed orbit in each fibre of $(B,A)\mapsto BA$}
For $\GL_r$ acting on factor pairs, fix $W$ of rank $k\le r$. Among the pairs with $BA=W$, those with $\rk B=\rk A=k$
form a single orbit. It is closed, and it lies in the closure of every orbit in the fibre $\{BA=W\}$. Consequently a
continuous $\GL_r$-invariant function is constant on each fibre of $(B,A)\mapsto BA$. For a complex reductive group
acting on a vector space, every fibre of the quotient map likewise contains exactly one closed orbit
\citep[Ch.~1]{xmumford1994geometric}.
\end{bgfact}

For $m=n=r=1$ (\bgref{D.30}), the fibre $\{ba=0\}$ consists of the origin and the two punctured axes, and the origin,
its only closed orbit, lies in the closure of both axes. This is the step behind \thmref{Proposition}{V.4}(a). A
continuous rule for combining LoRA factors that is invariant under each adapter's $\GL_r$ is automatically a function of
the adapters $B_iA_i$, also off the full-rank locus. Continuity matters. The projector onto $\operatorname{row}A$,
defined where $\rk A=r$, is invariant, yet it differs at $(0,A)$ and $(0,A')$, which have the same product $0$.

\fbAusedin{\thmref{Proposition}{V.4}, combining (\S\ref{sec:base-combine}).}

% ----------------------------------------------------------------------------------------------------------------------
\subsubsection*{Matrix and convexity tools}

\begin{bgfact}{D.32}{Ky Fan's maximum principle}
Let $S\in\R^{n\times n}$ be symmetric with eigenvalues $\lambda_1\ge\dots\ge\lambda_n$, and let $1\le k\le n$. Then
\[\max\{\tr(PSP\T):P\in\R^{k\times n},\ PP\T=I_k\}=\lambda_1+\dots+\lambda_k,\]
attained when the rows of $P$ span a top-$k$ eigenspace. For any matrix $X$ with singular values
$\sigma_1\ge\sigma_2\ge\cdots$, applying this to $S=X\T X$ gives $\max_P\|XP\T\|_F^2=\sigma_1^2+\dots+\sigma_k^2$, and
for square $X$ also $\|PXP\T\|_F^2\le\|XP\T\|_F^2\le\sigma_1^2+\dots+\sigma_k^2$. See \citet{xfan1949theorem} and
\citet[Ch.~III]{xbhatia1997matrix}.
\end{bgfact}

Among $A_0$ with $A_0A_0\T=cI$, LoRA's initial rate $\|GA_0\T\|_F^2$ under gradient flow is largest when
$\operatorname{row}A_0$ is the top-$r$ right singular space of $G$, which is LoRA-GA's \citep{wang2024lorab} choice for
$A_0$. For LOFT, the squared norm of the pulled-back gradient $P_r\skewop(W_0\T G)P_r\T$ is at most $2\sum_{k\le\lfloor
r/2\rfloor}\mu_k^2$, where $\pm i\mu_k$ are the eigenvalues of the skew matrix $\skewop(W_0\T G)$, whose singular values
are $\mu_1,\mu_1,\mu_2,\mu_2,\dots$. The floor appears because a skew matrix compressed to an odd-dimensional subspace
has a zero eigenvalue, and Cauchy interlacing bounds the rest.

\fbAusedin{LOFT (\S\ref{sec:erlangen-loft}), \thmref{Proposition}{II.3}.}

\begin{bgfact}{D.33}{Sylvester equations}
For $S\in\R^{p\times p}$, $T\in\R^{q\times q}$ and $C\in\R^{p\times q}$, the \textbf{Sylvester equation} $SX+XT=C$ has a
unique solution $X$ for every $C$ iff $S$ and $-T$ have no eigenvalue in common, since the eigenvalues of $X\mapsto
SX+XT$ are the sums $\lambda_i(S)+\lambda_j(T)$. If $S$ and $T$ are symmetric positive definite, this holds and
$X=\int_0^\infty e^{-tS}Ce^{-tT}\,dt$. The Bartels--Stewart algorithm \citep{xbartels1972algorithm} solves the equation
through Schur forms of $S$ and $T$ \citep[Ch.~4]{xhorn1991topics}.
\end{bgfact}

LoRA-Pro's equivalent gradient leaves a free $r\times r$ matrix $X$. The published method picks it by solving $B\T
B\,X+X\,AA\T=C$, with $C$ built from the gradient. On the full-rank locus both Gram matrices are positive definite, so
$X$ exists and is unique. The first-order weight velocity does not depend on $X$, while the resulting parameter update
is only $\Ort(r)$-equivariant.

\fbAusedin{\thmref{Theorem}{III.12}, when a change of chart does not matter (\S\ref{sec:fibres-naturality}),
  Figure~\ref{fig:gauge}.}

\begin{bgentry}{D.34}{the perspective map; convex combinations}
The \textbf{perspective map} is $\pi(Z,S)=S/Z$ on $\{(Z,S)\in\R\times\R^d:Z>0\}$. It is constant on rays,
$\pi(tZ,tS)=\pi(Z,S)$ for $t>0$, and it turns sums into convex combinations,
\[\pi(Z_1+Z_2,S_1+S_2)=\lambda\,\pi(Z_1,S_1)+(1-\lambda)\,\pi(Z_2,S_2),\qquad\lambda=\frac{Z_1}{Z_1+Z_2}.\]
\end{bgentry}

For a fixed query $q$, attention is $\mathrm{Attn}_q(K)=\pi(\Sigma_q(K))$, where
\[\Sigma_q(K)=\Big(\sum e^{\langle q,k\rangle/\sqrt d},\ \sum e^{\langle q,k\rangle/\sqrt d}\,v\Big)\]
is a monoid homomorphism on multisets of key--value pairs (\bgref{D.10}). With a prefix $P$ before the content $C$, the
identity above gives
\[\mathrm{Attn}_q(P\uplus C)=(1-\lambda_q)\,\mathrm{Attn}_q(C)+\lambda_q\,\mathrm{Attn}_q(P),\]
with $\mathrm{Attn}_q(P)$ in the convex hull $\operatorname{conv}V_P$ of the prefix values, so a prefix acts as a
parallel branch gated by its share of the mass.

\fbAusedin{\thmref{Theorem}{VI.5}, attention (\S\ref{sec:merge-attention}), Figure~\ref{fig:prefix}. \textsf{Reading:}
  \citet[\S 2.3.3 and \S 3.2.6]{xboyd2004convex}.}

